\documentclass[11pt]{article}

\usepackage[a4paper,margin=27mm,headheight=14pt]{geometry}
\usepackage[T1]{fontenc}
\usepackage[utf8]{inputenc}
\usepackage{lmodern}
\usepackage{microtype}
\usepackage{amsmath,amssymb,amsthm,mathtools,bm}
\allowdisplaybreaks
\usepackage{booktabs,longtable,array,tabularx,multirow}
\usepackage{enumitem}
\usepackage{xcolor}
\usepackage{hyperref}
\usepackage[nameinlink,capitalise,noabbrev]{cleveref}
\usepackage{tikz}
\usepackage{pgfplots}
\pgfplotsset{compat=1.18}
\usepackage{listings}
\usepackage{fancyhdr}
\usepackage{caption}
\usepackage{subcaption}
\usepackage{xurl}

\definecolor{frontierblue}{RGB}{22,73,118}
\definecolor{frontiergreen}{RGB}{25,110,83}
\definecolor{frontierorange}{RGB}{165,85,15}
\definecolor{softblue}{RGB}{237,245,252}
\definecolor{softgreen}{RGB}{237,248,243}
\definecolor{softorange}{RGB}{253,246,235}
\definecolor{codegray}{RGB}{246,246,246}

\hypersetup{
  pdftitle={Tail Quadrature, Exact Dyadic Error Laws, and q-Binomial Acceleration for the Fabius--Rvachev System},
  pdfauthor={OpenAI, based on the ProveIt corpus by Vladimir Reshetnikov},
  pdfsubject={New frontier results on the Fabius function, Rvachev up-function, Thue--Morse splines, quadrature, q-binomial acceleration, and infinite sinc products},
  pdfkeywords={Fabius function, Rvachev up-function, Thue-Morse, Gaussian quadrature, Gauss-Lobatto quadrature, q-binomial, sinc product, Peano kernel},
  colorlinks=true,
  linkcolor=frontierblue,
  citecolor=frontiergreen,
  urlcolor=frontierblue
}

\pagestyle{fancy}
\fancyhf{}
\fancyhead[L]{Fabius--Rvachev frontier report}
\fancyhead[R]{\thepage}
\renewcommand{\headrulewidth}{0.3pt}

\numberwithin{equation}{section}
\setcounter{tocdepth}{2}
\setcounter{secnumdepth}{3}
\setlist{itemsep=2pt,topsep=5pt}

\newtheorem{theorem}{Theorem}[section]
\newtheorem{proposition}[theorem]{Proposition}
\newtheorem{lemma}[theorem]{Lemma}
\newtheorem{corollary}[theorem]{Corollary}
\newtheorem{claim}[theorem]{Claim}
\theoremstyle{definition}
\newtheorem{definition}[theorem]{Definition}
\newtheorem{algorithm}[theorem]{Algorithm}
\newtheorem{example}[theorem]{Example}
\theoremstyle{remark}
\newtheorem{remark}[theorem]{Remark}
\newtheorem{warning}[theorem]{Warning}
\newtheorem{status}[theorem]{Status note}

\DeclareMathOperator{\supp}{supp}
\DeclareMathOperator{\Var}{Var}
\DeclareMathOperator{\Cov}{Cov}
\DeclareMathOperator{\diag}{diag}
\DeclareMathOperator{\sgn}{sgn}
\DeclareMathOperator{\spanop}{span}
\newcommand{\R}{\mathbb R}
\newcommand{\Q}{\mathbb Q}
\newcommand{\Z}{\mathbb Z}
\newcommand{\N}{\mathbb N}
\newcommand{\C}{\mathbb C}
\newcommand{\E}{\mathbb E}
\newcommand{\Prob}{\mathbb P}
\newcommand{\Up}{\operatorname{up}}
\newcommand{\F}{F}
\newcommand{\eps}{\varepsilon}
\newcommand{\sinc}{\operatorname{sinc}}
\newcommand{\dd}{\,\mathrm d}
\newcommand{\e}{\mathrm e}
\newcommand{\ii}{\mathrm i}
\newcommand{\Oh}{\mathrm O}
\newcommand{\littleoh}{\mathrm o}
\newcommand{\qbinom}[3]{\genfrac{[}{]}{0pt}{}{#1}{#2}_{#3}}
\newcommand{\abs}[1]{\left|#1\right|}
\newcommand{\norm}[1]{\left\lVert#1\right\rVert}
\newcommand{\set}[1]{\left\{#1\right\}}
\newcommand{\floor}[1]{\left\lfloor#1\right\rfloor}
\newcommand{\ceil}[1]{\left\lceil#1\right\rceil}
\newcommand{\angles}[1]{\left\langle#1\right\rangle}
\newcommand{\pospart}[1]{\left(#1\right)_{+}}
\newcommand{\aScale}{a_n}
\newcommand{\TM}{\varepsilon}

\newcommand{\frontierbox}[1]{%
\begin{center}
\fcolorbox{frontierblue}{softblue}{\parbox{0.91\textwidth}{#1}}
\end{center}}
\newcommand{\statusbox}[1]{%
\begin{center}
\fcolorbox{frontierorange}{softorange}{\parbox{0.91\textwidth}{#1}}
\end{center}}

\lstdefinestyle{code}{
  backgroundcolor=\color{codegray},
  basicstyle=\ttfamily\small,
  columns=fullflexible,
  keepspaces=true,
  frame=single,
  rulecolor=\color{black!25},
  breaklines=true,
  showstringspaces=false,
  keywordstyle=\color{frontierblue}\bfseries,
  commentstyle=\color{frontiergreen},
  stringstyle=\color{frontierorange}
}

\title{\textbf{Tail Quadrature, Exact Dyadic Error Laws,\\
and $q$-Binomial Acceleration for the Fabius--Rvachev System}\\[0.65em]
\large New frontier results connecting Thue--Morse splines, positive Peano kernels, orthogonal polynomials, and infinite sinc products}
\author{A research synthesis based on the \texttt{ProveIt} Fabius-function corpus}
\date{27 August 2026}

\begin{document}
\maketitle

\begin{abstract}
The Fabius function and Rvachev's compactly supported up-function sit at an unusually fertile intersection of probability, dyadic self-similarity, Thue--Morse cancellation, exact rational arithmetic, $q$-binomial interpolation, and lacunary Fourier analysis.  The accompanying \texttt{ProveIt} corpus already contains a broad human-readable exposition, a Lean walkthrough, a mathematical glossary, exact dyadic evaluators, a bridge from coefficient formulas to endpoint asymptotics, sharp repeated-integration laws, extensive Thue--Morse formula atlases, and a detailed analysis of the infinite sinc product.  The purpose of the present report is therefore not to repeat that material, but to isolate gaps explicitly left by the corpus and to push them to theorem-level conclusions.

The first new result is an exact uniform and $C^m$ error law for the \emph{ordinary finite convolution prefix}
\[
 p_n=\beta_{2^{-1}}*\beta_{2^{-2}}*\cdots*\beta_{2^{-n}},
\]
a case that the existing repeated-integration report deliberately left open.  For every $m\ge0$ and $n\ge m+3$,
\[
 \norm{p_n^{(m)}-\Up^{(m)}}_\infty
 =\frac{2^{\binom{m+3}{2}}}{18}\,4^{-n},
\]
and equality occurs at the explicit dyadic point $-1+2^{-(m+2)}$.  In particular,
\(\norm{p_n-\Up}_\infty=\frac49 4^{-n}\) for $n\ge3$.

The second result resolves a more ambitious direction suggested in the corpus: replace the exact self-similar tail by a \emph{positive compactly supported law matching additional moments}.  Gaussian quadrature for the up-law gives an $r$-node positive tail model with exact rate $4^{-rn}$, while a Gauss--Lobatto variant includes the endpoints and therefore preserves the support $[-1,1]$ exactly.  Both families admit sharp $C^m$ constants, explicit extremal points, nonnegative Peano kernels, and rational constants expressed by Hankel determinants of the up moments.  The first support-preserving variance-matched scheme is
\[
 \frac1{18}p_n(x+2^{-n})+\frac89p_n(x)+\frac1{18}p_n(x-2^{-n}),
\]
whose exact uniform error is \(\frac{7168}{2025}16^{-n}\) for $n\ge5$.

Finally, the reciprocal infinite sinc product yields an all-orders expansion of the ordinary prefix, and the geometric error scale $q=1/4$ produces a new Richardson family whose weights are Gaussian $q$-binomial coefficients.  These results supply a direct four-way bridge:
\[
\begin{gathered}
\text{Thue--Morse plateaux}
\longleftrightarrow
\text{positive quadrature kernels}\\
\longleftrightarrow
\text{reciprocal sinc coefficients}
\longleftrightarrow
q\text{-binomial acceleration}.
\end{gathered}
\]
All proofs are given in ordinary analysis, with exact rational and high-precision checks.  The results are new relative to the audited corpus.  The finite geometric \(q\)-binomial/Lagrange bridge identified below now has matching Lean declarations; the analytic approximation, error, remainder, conditioning, and convergence theorems are not thereby claimed formalized or externally peer reviewed.
\end{abstract}

\tableofcontents
\newpage

\section{Corpus audit, novelty boundary, and theorem status}
\label{sec:audit}

\subsection{What was audited}

The source tree under \path{Analysis/FabiusFunction/docs} has several layers rather than a single linear manuscript.  The audit used the current document families as follows.

\begin{center}
\begin{tabularx}{0.97\textwidth}{@{}p{0.31\textwidth}X@{}}
\toprule
Document family & Mathematical role in the audit\\
\midrule
\path{Fabius_Function_and_Rvachev_Up} & Canonical definitions, normalizations, probability law, infinite convolution and sinc product, exact dyadic arithmetic, $q$-binomial/Thue--Morse formulas, and endpoint asymptotics.\\
\path{FabiusFunction_Mathematical_Glossary} & Terminology and coverage map across analysis, probability, Fourier methods, $q$-calculus, computability, and formalization.\\
\path{fabius_lean_walkthrough} & Structure and dependency map of the Lean development; used to separate established formal results from exploratory claims.\\
\path{Non_Elementarity_of_the_Fabius_Function} & Differential-algebraic and non-elementarity direction; audited to avoid misclassifying an approximation theorem as a transcendence result.\\
Exact-value and asymptotic papers & Dyadic evaluators, alternative representations, repeated integration, Thue--Morse convergence, and the coefficient-to-saddle bridge.\\
Canonical non-formalized frontier volume & Consolidated research claims, open directions, and explicit formalization boundary.  Its absorbed notebooks were treated as one mathematical corpus rather than as independent evidence.\\
Active drafts on Thue--Morse formulas and Fourier decay & Recent formula atlases and sharp analysis of the infinite sinc product; used to avoid duplicating the newest work.\\
Archive and audit material & Provenance, superseded variants, and contradiction checks.  Byte-identical or explicitly frozen historical copies were not counted as separate mathematical contributions.\\
\bottomrule
\end{tabularx}
\end{center}

The principal baseline established by this corpus includes:
\begin{enumerate}
\item the random-series and infinite-convolution construction of $\Up$;
\item the exact Fourier product \(\widehat\Up(\xi)=\prod_{k\ge1}\sinc(2^{-k}\xi)\);
\item exact rational values at dyadic arguments and executable recurrences;
\item finite Thue--Morse and Gaussian $q$-binomial formulas;
\item sharp endpoint asymptotics with periodic corrections;
\item a sharp law for the repeated-integration approximation, whose last dyadic box is duplicated;
\item extensive typical-ray and exceptional-ray results for the Fourier product.
\end{enumerate}

\subsection{The two explicit gaps used as starting points}

The repeated-integration report compares its support-corrected spline $f_n$ with the ordinary prefix $p_n$.  It proves
\[
 \norm{f_n-\Up}_\infty=\frac{8}{9\,4^n}
\]
from the second stage onward, but labels the corresponding exact law for $p_n$ as outside the scope of that paper.  It then observes that a positive compactly supported tail model matching the true variance would cancel the leading second-derivative term and should produce a $16^{-n}$ scheme, while leaving positivity and support as nontrivial constraints.

Those two sentences are the main launch point of this report.  The first becomes \cref{thm:ordinary-exact}; the second becomes the Gaussian and Gauss--Lobatto hierarchies in \cref{sec:gaussian,sec:lobatto}.  The Gauss--Lobatto construction is especially decisive: it is positive, matches the variance, and includes the two endpoints, so the final spline has support exactly $[-1,1]$.

\subsection{Status discipline}

\statusbox{\textbf{Mathematical status.}  Every displayed frontier theorem below is accompanied by a proof or by an explicit reduction to standard Gaussian/Lobatto quadrature facts.  The low-order constants were independently checked by exact rational arithmetic or high-precision evaluation.  The results are presented as new relative to the audited \texttt{ProveIt} corpus.  A broad literature search did not locate these Fabius-specific formulas, but this report does not assert an exhaustive priority claim.  Only the finite algebraic bridge explicitly cross-referenced to matching declarations in \path{GeometricQBinomialLagrange.lean} is labelled Lean-formalized; the report's analytic theorem families remain frontier.}

\subsection{Executive list of new results}

\frontierbox{\textbf{New results proved in this report.}
\begin{enumerate}[leftmargin=2em]
\item An exact ordinary-prefix law in every fixed $C^m$ norm, including exact extremal dyadic points.
\item A positive $r$-node Gaussian tail hierarchy with sharp rate $4^{-rn}$ and exact constants \(h_r/(2r)!\).
\item A positive support-exact Gauss--Lobatto hierarchy with the same rate and constants \(\ell_r/(2r)!\).
\item A two-sided order for all $2r$-convex observables, placing the true tail between its Gaussian and Lobatto surrogates.
\item Rational Hankel-determinant formulas for all sharp constants, tied to the rational moment sequence of the infinite sinc product.
\item A complete reciprocal-sinc expansion of the ordinary prefix in powers of $4^{-n}$.
\item A $q=1/4$ Richardson family whose weights are explicit Gaussian $q$-binomial coefficients.
\end{enumerate}}

\section{The master dictionary: probability, sinc products, and Thue--Morse splines}
\label{sec:dictionary}

\subsection{The up-law and its finite prefixes}

For $a>0$, write
\begin{equation}
 \beta_a(x):=\frac1{2a}\mathbf 1_{[-a,a]}(x).
 \label{eq:box}
\end{equation}
Let $U_1,U_2,\ldots$ be independent random variables uniformly distributed on $[-1,1]$, and define
\begin{equation}
 Y:=\sum_{k=1}^{\infty}2^{-k}U_k,
 \qquad
 S_n:=\sum_{k=1}^{n}2^{-k}U_k.
 \label{eq:Y-Sn}
\end{equation}
The series converges absolutely, $Y\in[-1,1]$, and its density is Rvachev's even up-function.  We write
\begin{equation}
 \mu:=\mathcal L(Y),
 \qquad
 p_n:=\beta_{2^{-1}}*\beta_{2^{-2}}*\cdots*\beta_{2^{-n}}.
 \label{eq:mu-pn}
\end{equation}
Thus $p_n$ is the density of $S_n$ and
\begin{equation}
 \supp(p_n)=[-1+2^{-n},1-2^{-n}].
 \label{eq:pn-support}
\end{equation}
The omitted tail is an exact scaled copy of the original law:
\begin{equation}
 Y-S_n\ \stackrel{d}=\ 2^{-n}Y.
 \label{eq:tail-self-similar}
\end{equation}
If $\mu_a$ denotes the pushforward of $\mu$ under $x\mapsto ax$, then, with $a=2^{-n}$,
\begin{equation}
 \boxed{\Up=p_n*\mu_a.}
 \label{eq:master-tail-factorization}
\end{equation}
This elementary identity is the source of every approximation theorem below.

We use the Fourier convention
\begin{equation}
 \widehat g(\xi)=\int_{\R}g(x)\e^{-\ii x\xi}\dd x,
 \qquad
 \sinc z=\frac{\sin z}{z},\quad \sinc0=1.
 \label{eq:Fourier-convention}
\end{equation}
Then
\begin{align}
 P_n(\xi):=\widehat p_n(\xi)
  &=\prod_{k=1}^{n}\sinc(2^{-k}\xi),\label{eq:Pn}\\
 \Phi(\xi):=\widehat\Up(\xi)
  &=\prod_{k=1}^{\infty}\sinc(2^{-k}\xi),\label{eq:Phi}\\
 \Phi(\xi)&=P_n(\xi)\Phi(2^{-n}\xi).\label{eq:Phi-factor}
\end{align}

On the unit interval, the bounded Fabius function is the shifted left half of the bump:
\begin{equation}
 F(x)=\Up(x-1),\qquad 0\le x\le1.
 \label{eq:F-shift}
\end{equation}
Consequently every approximation to $\Up$ immediately gives one for $F$ by translating the argument by one.

\subsection{The finite Thue--Morse formula}

Let
\begin{equation}
 \TM_j:=(-1)^{s_2(j)},
 \label{eq:TM-sign}
\end{equation}
where $s_2(j)$ is the binary digit sum.  The generalized Irwin--Hall formula for unequal dyadic boxes becomes a Thue--Morse spline.

\begin{proposition}[Finite Thue--Morse box-spline formula]
\label{prop:TM-spline}
For $n\ge1$,
\begin{equation}
 \boxed{
 p_n(x)=\frac{2^{\binom n2}}{(n-1)!}
 \sum_{j=0}^{2^n-1}\TM_j
 \pospart{x+1-2^{-n}-j2^{1-n}}^{n-1}.}
 \label{eq:TM-spline}
\end{equation}
For $0\le r\le n-1$, differentiation gives
\begin{equation}
 p_n^{(r)}(x)=\frac{2^{\binom n2}}{(n-1-r)!}
 \sum_{j=0}^{2^n-1}\TM_j
 \pospart{x+1-2^{-n}-j2^{1-n}}^{n-1-r}
 \label{eq:TM-spline-derivative}
\end{equation}
away from the knots, and distributionally everywhere.
\end{proposition}

\begin{proof}
Write $2^{-k}U_k=V_k-2^{-k}$ with $V_k$ uniform on $[0,2^{1-k}]$.  The density of $\sum V_k$ is the inclusion--exclusion spline
\[
 \frac1{(n-1)!\prod_{k=1}^n2^{1-k}}
 \sum_{A\subseteq\{1,\ldots,n\}}(-1)^{|A|}
 \left(t-\sum_{k\in A}2^{1-k}\right)_+^{n-1}.
\]
Now \(\prod_{k=1}^n2^{1-k}=2^{-\binom n2}\), the centering shift is
\(\sum_{k=1}^n2^{-k}=1-2^{-n}\), and the map
\[
 A\longmapsto j(A):=\sum_{k\in A}2^{n-k}
\]
is a bijection from subsets to $\{0,\ldots,2^n-1\}$.  It satisfies
\(\sum_{k\in A}2^{1-k}=j(A)2^{1-n}\) and
\((-1)^{|A|}=(-1)^{s_2(j(A))}\), proving \eqref{eq:TM-spline}.  Differentiation gives \eqref{eq:TM-spline-derivative}.
\end{proof}

\subsection{A saturated derivative plateau}

The sharp constants below come from a simple but powerful fact about Thue--Morse prefix sums:
\begin{equation}
 \sum_{j=0}^{2q-1}\TM_j=0,
 \qquad
 \sum_{j=0}^{2q}\TM_j=\TM_q.
 \label{eq:TM-prefix-two-valued}
\end{equation}
Thus every partial sum is $0$ or $\pm1$.

\begin{lemma}[Thue--Morse derivative plateau]
\label{lem:plateau}
For $r\ge0$ put
\begin{equation}
 C_r:=2^{\binom{r+1}{2}},
 \qquad
 x_r:=-1+2^{-r}.
 \label{eq:Cr-xr}
\end{equation}
If $n\ge r+1$, then
\begin{equation}
 \boxed{\norm{p_n^{(r)}}_\infty=C_r.}
 \label{eq:derivative-norm}
\end{equation}
Moreover,
\begin{equation}
 p_n^{(r)}(x)=C_r
 \quad\text{for almost every }x\in[x_r-2^{-n},x_r+2^{-n}].
 \label{eq:plateau}
\end{equation}
\end{lemma}

\begin{proof}
First take $n=r+1$.  In \eqref{eq:TM-spline-derivative}, the $r$th derivative is
\[
 p_{r+1}^{(r)}(x)=C_r
 \sum_{\substack{0\le j<2^{r+1}\\ j2^{-r}<x+1-2^{-(r+1)}}}\TM_j.
\]
By \eqref{eq:TM-prefix-two-valued}, its absolute value never exceeds $C_r$.  On the first knot interval
\[
 [-1+2^{-(r+1)},\,-1+3\cdot2^{-(r+1)}],
\]
only the $j=0$ term is active, so the value is exactly $C_r$.  This interval has center $x_r$ and half-width $2^{-(r+1)}$.

For $n>r+1$, factor
\[
 p_n=p_{r+1}*\beta_{2^{-(r+2)}}*\cdots*\beta_{2^{-n}}.
\]
The remaining convolution factor is a probability density supported on an interval of radius
\[
 \sum_{k=r+2}^{n}2^{-k}=2^{-(r+1)}-2^{-n}.
\]
Convolution with a probability density cannot increase the $L^\infty$ norm.  It preserves the constant value $C_r$ wherever the translated support remains inside the first plateau, leaving precisely the interval in \eqref{eq:plateau}.  Hence the upper bound is attained.
\end{proof}

\begin{remark}
The same binary fact controls three phenomena already prominent in the corpus: the finite spline formula, the exact derivative plateau, and the terminating derivative tower of $F$ at dyadic arguments.  In the present report it acquires a fourth role: it turns every positive Peano-kernel upper bound into an exact equality.
\end{remark}

\section{Rational moments from the infinite sinc product}
\label{sec:moments}

\subsection{Moment generating function and cumulants}

The moment generating function of $Y$ is entire and factors as
\begin{equation}
 M(t):=\E\e^{tY}
 =\prod_{k=1}^{\infty}\frac{\sinh(t/2^k)}{t/2^k}.
 \label{eq:MGF}
\end{equation}
The classical expansion
\begin{equation}
 \log\frac{\sinh z}{z}
 =\sum_{m=1}^{\infty}
 \frac{2^{2m-1}B_{2m}}{m(2m)!}z^{2m}
 \label{eq:logsinh}
\end{equation}
therefore yields
\begin{equation}
 \boxed{
 \log M(t)=\sum_{m=1}^{\infty}
 \frac{2^{2m-1}B_{2m}}{m(2m)!(4^m-1)}t^{2m}.}
 \label{eq:logM}
\end{equation}
If $\kappa_j$ denotes the $j$th cumulant, then
\begin{equation}
 \boxed{
 \kappa_{2m}=\frac{2^{2m-1}B_{2m}}{m(4^m-1)},
 \qquad
 \kappa_{2m+1}=0.}
 \label{eq:cumulants}
\end{equation}
The factors $4^m-1$ are the moment-side trace of the dyadic sinc scales.

Let $\mu_j:=\E Y^j$.  The standard moment--cumulant recurrence gives
\begin{equation}
 \mu_0=1,
 \qquad
 \mu_n=\sum_{k=1}^{n}\binom{n-1}{k-1}\kappa_k\mu_{n-k}.
 \label{eq:moment-recurrence}
\end{equation}
In particular, every moment is rational.  The first values needed below are
\begin{align}
 \mu_0&=1,&
 \mu_2&=\frac19,&
 \mu_4&=\frac{19}{675},&
 \mu_6&=\frac{583}{59535},\label{eq:moments-low}\\
 \mu_8&=\frac{132809}{32531625},&
 \mu_{10}&=\frac{46840699}{24405225075},&
 \mu_{12}&=\frac{4068990560161}{4133856862760625}.\nonumber
\end{align}
All odd moments vanish.

\subsection{Hankel determinants and exact quadrature constants}

Define the Hankel determinants
\begin{equation}
 \Delta_r:=\det(\mu_{i+j})_{0\le i,j\le r-1},
 \qquad \Delta_0:=1.
 \label{eq:Hankel}
\end{equation}
Because $\mu$ has a positive smooth density on $(-1,1)$, every $\Delta_r$ is positive.  Let $\pi_r$ be the monic orthogonal polynomial of degree $r$ for $\mu$ and put
\begin{equation}
 h_r:=\int_{-1}^{1}\pi_r(x)^2\dd\mu(x).
 \label{eq:hr}
\end{equation}
Then the Gram-determinant identity gives
\begin{equation}
 \boxed{h_r=\frac{\Delta_{r+1}}{\Delta_r}.}
 \label{eq:hr-determinant}
\end{equation}
Thus every $h_r$ is a positive rational number.  Equivalently,
\begin{equation}
 h_r=\min\left\{\int P(x)^2\dd\mu(x):P\text{ monic of degree }r\right\}.
 \label{eq:hr-minimizer}
\end{equation}

For the Lobatto construction, introduce the weighted moment functional
\begin{equation}
 \int q(x)\dd\nu(x):=\int(1-x^2)q(x)\dd\mu(x)
 \label{eq:nu-weighted}
\end{equation}
and its Hankel determinants
\begin{equation}
 \Lambda_s:=\det(\mu_{i+j}-\mu_{i+j+2})_{0\le i,j\le s-1},
 \qquad \Lambda_0:=1.
 \label{eq:Lambda}
\end{equation}
If $\chi_{r-1}$ is monic of degree $r-1$ and orthogonal for $\nu$, define
\begin{equation}
 \ell_r:=\int(1-x^2)\chi_{r-1}(x)^2\dd\mu(x).
 \label{eq:ellr}
\end{equation}
Then
\begin{equation}
 \boxed{\ell_r=\frac{\Lambda_r}{\Lambda_{r-1}}}
 \label{eq:ellr-determinant}
\end{equation}
and
\begin{equation}
 \ell_r=\min\left\{\int(1-x^2)P(x)^2\dd\mu(x):P\text{ monic of degree }r-1\right\}.
 \label{eq:ellr-minimizer}
\end{equation}
These rational norms will be the exact sharp constants of the two quadrature hierarchies.

\subsection{Low-degree orthogonal data}

Symmetry forces the expected parity.  From \eqref{eq:moments-low},
\begin{align}
 \pi_1(x)&=x,& h_1&=\frac19,\label{eq:pi1}\\
 \pi_2(x)&=x^2-\frac19,& h_2&=\frac{32}{2025},\label{eq:pi2}\\
 \pi_3(x)&=x^3-\frac{19}{75}x,& h_3&=\frac{19808}{7441875}.
 \label{eq:pi3}
\end{align}
For the weighted Lobatto measure,
\begin{align}
 \chi_1(x)&=x,& \ell_2&=\frac{56}{675},\label{eq:chi1}\\
 \chi_2(x)&=x^2-\frac7{75},& \ell_3&=\frac{78976}{7441875}.
 \label{eq:chi2}
\end{align}
The rational coefficients coexist with algebraic quadrature nodes; this arithmetic split is one of the useful structural features of the construction.

\section{The exact ordinary-prefix error law}
\label{sec:ordinary}

The ordinary prefix simply omits the tail.  At first sight this seems less structured than the support-corrected repeated-integration scheme.  In fact the omitted tail has a canonical positive second-order potential, and the same Thue--Morse plateau used in the repeated-integration analysis makes the resulting bound exact.

\subsection{A convex-order potential for the omitted tail}

Define
\begin{equation}
 J(t):=\frac12\left(\E\abs{t-Y}-\abs t\right).
 \label{eq:J-def}
\end{equation}

\begin{lemma}[Positive tail potential]
\label{lem:J}
The function $J$ is even, nonnegative, continuous, and supported on $[-1,1]$.  In the sense of distributions,
\begin{equation}
 J''=\mu-\delta_0.
 \label{eq:J-second}
\end{equation}
Moreover,
\begin{equation}
 \int_{\R}J(t)\dd t=\frac{\mu_2}{2}=\frac1{18}.
 \label{eq:J-mass}
\end{equation}
\end{lemma}

\begin{proof}
The law of $Y$ is centered.  Jensen's inequality for the convex map $y\mapsto\abs{t-y}$ gives
\[
 \E\abs{t-Y}\ge\abs{t-\E Y}=\abs t,
\]
so $J\ge0$.  If $t\ge1$, then $t-Y\ge0$ and
\(\E\abs{t-Y}=t-\E Y=t\); the case $t\le-1$ is symmetric.  Hence $J$ vanishes outside $[-1,1]$.

Since \((\abs{t-y}/2)''=\delta_y\), averaging in $y$ and subtracting the case $y=0$ yields \eqref{eq:J-second}.  Finally, $J$ and its first derivative vanish at the boundary of its support, so two integrations by parts give
\[
 2\int J(t)\dd t
 =\int t^2J''(t)\dd t
 =\int t^2\dd\mu(t)-0
 =\mu_2=\frac19.
\]
\end{proof}

For $a>0$, set
\begin{equation}
 J_a(t):=aJ(t/a).
 \label{eq:Ja}
\end{equation}
Then
\begin{equation}
 J_a''=\mu_a-\delta_0,
 \qquad
 \supp(J_a)\subseteq[-a,a],
 \qquad
 \int J_a=\frac{a^2}{18}.
 \label{eq:Ja-properties}
\end{equation}
Combining this with \eqref{eq:master-tail-factorization} gives the exact factorization
\begin{equation}
 \boxed{\Up-p_n=p_n*(\mu_a-\delta_0)=p_n''*J_a,\qquad a=2^{-n}.}
 \label{eq:ordinary-factorization}
\end{equation}

\subsection{Sharp $C^m$ constants}

\begin{theorem}[Exact ordinary-prefix law]
\label{thm:ordinary-exact}
Let $m\ge0$ and $n\ge m+3$.  Then
\begin{equation}
 \boxed{
 \norm{p_n^{(m)}-\Up^{(m)}}_\infty
 =\frac{2^{\binom{m+3}{2}}}{18}\,4^{-n}.}
 \label{eq:ordinary-exact}
\end{equation}
At the dyadic point
\begin{equation}
 x_{m+2}=-1+2^{-(m+2)},
 \label{eq:ordinary-extremal-point}
\end{equation}
we have the signed identity
\begin{equation}
 \boxed{
 \bigl(p_n^{(m)}-\Up^{(m)}\bigr)(x_{m+2})
 =-\frac{2^{\binom{m+3}{2}}}{18}\,4^{-n}.}
 \label{eq:ordinary-signed}
\end{equation}
\end{theorem}

\begin{proof}
Differentiate \eqref{eq:ordinary-factorization} $m$ times:
\begin{equation}
 \Up^{(m)}-p_n^{(m)}=p_n^{(m+2)}*J_a.
 \label{eq:ordinary-m-factorization}
\end{equation}
The kernel is nonnegative and has mass $a^2/18$.  By \cref{lem:plateau},
\[
 \norm{p_n^{(m+2)}}_\infty=C_{m+2}=2^{\binom{m+3}{2}}.
\]
Therefore
\[
 \norm{\Up^{(m)}-p_n^{(m)}}_\infty
 \le C_{m+2}\frac{a^2}{18}.
\]
At $x=x_{m+2}$, the whole integration interval
\([x-a,x+a]\) lies inside the plateau from \eqref{eq:plateau}, because $n\ge m+3$.  Hence
\[
 (p_n^{(m+2)}*J_a)(x_{m+2})
 =C_{m+2}\int J_a
 =C_{m+2}\frac{a^2}{18}.
\]
Thus the upper bound is attained, and the sign follows from \eqref{eq:ordinary-m-factorization}.
\end{proof}

\begin{corollary}[Uniform law]
\label{cor:ordinary-uniform}
For every $n\ge3$,
\begin{equation}
 \boxed{\norm{p_n-\Up}_\infty=\frac49\,4^{-n}.}
 \label{eq:ordinary-uniform}
\end{equation}
The negative extremum occurs at $x=-3/4$:
\begin{equation}
 p_n(-3/4)-\Up(-3/4)=-\frac49\,4^{-n}.
 \label{eq:ordinary-at-34}
\end{equation}
\end{corollary}

\begin{corollary}[Bounded Fabius version]
\label{cor:F-ordinary}
Define $P_n(x):=p_n(x-1)$ on $[0,1]$.  Then, for $m\ge0$ and $n\ge m+3$,
\begin{equation}
 \norm{P_n^{(m)}-F^{(m)}}_{L^\infty[0,1]}
 =\frac{2^{\binom{m+3}{2}}}{18}\,4^{-n},
 \label{eq:F-ordinary}
\end{equation}
with negative equality at $x=2^{-(m+2)}$.
\end{corollary}

\begin{remark}[Repeated integration]
The support-corrected repeated-integration density is
\(f_n=p_n*\beta_{2^{-n}}\).  The corpus proves
\[
 \norm{f_n^{(m)}-\Up^{(m)}}_\infty
 =\frac{2^{\binom{m+3}{2}}}{9}\,4^{-n}
\]
for the stable range, exactly twice the ordinary-prefix constant and with the opposite local sign at the corresponding saturated plateau.  The two laws are governed by positive second-order kernels of masses $1/18$ and $1/9$ respectively.
\end{remark}

\subsection{Why the equality is finite-stage exact}

A formal Taylor expansion would predict
\[
 p_n=\Up-\frac1{18}4^{-n}\Up''+\Oh(16^{-n}).
\]
At $x=-3/4$, however, the theorem is not merely asymptotic.  The kernel $J_a$ is supported entirely inside a region where $p_n''$ is exactly the constant $8$, so all higher corrections collapse into the same finite convolution identity.  This is the same mechanism that made the repeated-integration error exact, but with the omitted-tail potential replacing the uniform-tail potential.

\section{Gaussian tail quadrature: a positive $4^{-rn}$ hierarchy}
\label{sec:gaussian}

The ordinary prefix corresponds to replacing the true tail by the single atom $\delta_0$.  Gaussian quadrature replaces it by $r$ positive atoms while matching the first $2r-1$ moments.  Because the tail scale is $a=2^{-n}$, matching through degree $2r-1$ forces the first possible error to have size $a^{2r}=4^{-rn}$.

\subsection{The Gaussian rule for the up-law}

Let $\pi_r$ be the monic orthogonal polynomial from \eqref{eq:hr}.  Its roots
\begin{equation}
 -1<\xi_{r,1}<\cdots<\xi_{r,r}<1
 \label{eq:gauss-nodes}
\end{equation}
are simple.  There are unique positive weights $w_{r,j}$ such that
\begin{equation}
 Q_r:=\sum_{j=1}^{r}w_{r,j}\delta_{\xi_{r,j}}
 \label{eq:Qr}
\end{equation}
is exact on polynomials of degree at most $2r-1$:
\begin{equation}
 \int q\dd Q_r=\int q\dd\mu,
 \qquad \deg q\le2r-1.
 \label{eq:Gauss-exactness}
\end{equation}
Because $\mu$ is even, the nodes and weights occur in symmetric pairs, with a central node at $0$ when $r$ is odd.

\begin{proposition}[Positive Gaussian Peano kernel]
\label{prop:gauss-kernel}
The error functional
\begin{equation}
 \mathcal E_r^G(\phi):=\int\phi\dd\mu-\int\phi\dd Q_r
 \label{eq:Gauss-functional}
\end{equation}
is nonnegative whenever $\phi\in C^{2r}[-1,1]$ and $\phi^{(2r)}\ge0$.  It has a nonnegative Peano kernel $K_r^G$, supported on $[-1,1]$, such that
\begin{equation}
 \mathcal E_r^G(\phi)=\int_{-1}^{1}\phi^{(2r)}(t)K_r^G(t)\dd t,
 \label{eq:Gauss-Peano}
\end{equation}
\begin{equation}
 (K_r^G)^{(2r)}=\mu-Q_r,
 \qquad
 \int K_r^G=\frac{h_r}{(2r)!}.
 \label{eq:Gauss-kernel-mass}
\end{equation}
\end{proposition}

\begin{proof}
Let $H$ be the Hermite interpolant of degree at most $2r-1$ satisfying
\[
 H(\xi_{r,j})=\phi(\xi_{r,j}),
 \qquad
 H'(\xi_{r,j})=\phi'(\xi_{r,j})
 \quad(1\le j\le r).
\]
The remainder formula is
\begin{equation}
 \phi(x)-H(x)=\frac{\phi^{(2r)}(\eta_x)}{(2r)!}\pi_r(x)^2
 \label{eq:Hermite-Gauss-remainder}
\end{equation}
for some $\eta_x\in(-1,1)$.  Hence $\phi-H\ge0$ when $\phi^{(2r)}\ge0$.  Gaussian exactness and interpolation give
\[
 \mathcal E_r^G(\phi)
 =\int(\phi-H)\dd\mu\ge0.
\]
The Peano-kernel theorem now gives a nonnegative kernel.  For $\phi(x)=x^{2r}$, the difference $\phi-H$ is the monic degree-$2r$ polynomial with double zeros at all Gaussian nodes, hence equals $\pi_r^2$.  Therefore
\[
 \mathcal E_r^G(x^{2r})=\int\pi_r^2\dd\mu=h_r.
\]
Since $(x^{2r})^{(2r)}=(2r)!$, the kernel mass is $h_r/(2r)!$.  The distributional derivative identity follows directly from \eqref{eq:Gauss-Peano}.
\end{proof}

For $a>0$, let $Q_{r,a}$ be the pushforward of $Q_r$ under $x\mapsto ax$, and set
\begin{equation}
 K_{r,a}^G(t):=a^{2r-1}K_r^G(t/a).
 \label{eq:scaled-G-kernel}
\end{equation}
Then
\begin{equation}
 (K_{r,a}^G)^{(2r)}=\mu_a-Q_{r,a},
 \qquad
 \int K_{r,a}^G=a^{2r}\frac{h_r}{(2r)!}.
 \label{eq:scaled-G-properties}
\end{equation}

\subsection{The Gaussian-corrected finite spline}

Define
\begin{equation}
 G_{n,r}:=p_n*Q_{r,2^{-n}}
 =\sum_{j=1}^{r}w_{r,j}\,p_n(\,\cdot-2^{-n}\xi_{r,j}).
 \label{eq:Gnr}
\end{equation}
This is a positive probability density and a finite mixture of shifted Thue--Morse splines.  Its support is strictly contained in $[-1,1]$ because all Gaussian nodes lie in $(-1,1)$.

\begin{theorem}[Sharp Gaussian hierarchy]
\label{thm:gaussian-exact}
Let $r\ge1$, $m\ge0$, and $n\ge m+2r+1$.  Then
\begin{equation}
 \boxed{
 \norm{G_{n,r}^{(m)}-\Up^{(m)}}_\infty
 =\frac{2^{\binom{m+2r+1}{2}}h_r}{(2r)!}\,2^{-2rn}.}
 \label{eq:Gaussian-exact}
\end{equation}
At
\begin{equation}
 x_{m+2r}=-1+2^{-(m+2r)},
 \label{eq:G-extremal}
\end{equation}
we have
\begin{equation}
 \boxed{
 \bigl(G_{n,r}^{(m)}-\Up^{(m)}\bigr)(x_{m+2r})
 =-\frac{2^{\binom{m+2r+1}{2}}h_r}{(2r)!}\,2^{-2rn}.}
 \label{eq:G-signed}
\end{equation}
\end{theorem}

\begin{proof}
Using \eqref{eq:master-tail-factorization} and \eqref{eq:scaled-G-properties}, with $a=2^{-n}$,
\begin{align}
 \Up-G_{n,r}
 &=p_n*(\mu_a-Q_{r,a})\nonumber\\
 &=p_n^{(2r)}*K_{r,a}^G.
 \label{eq:G-factorization}
\end{align}
After $m$ derivatives,
\[
 \Up^{(m)}-G_{n,r}^{(m)}
 =p_n^{(m+2r)}*K_{r,a}^G.
\]
The nonnegative kernel has mass $a^{2r}h_r/(2r)!$.  The derivative plateau lemma gives the sharp upper bound
\[
 C_{m+2r}\,a^{2r}\frac{h_r}{(2r)!}.
\]
At $x_{m+2r}$ the full kernel support $[-a,a]$ lies inside the constant positive plateau, so equality holds.  The sign follows from the displayed factorization.
\end{proof}

\begin{corollary}[Rate ladder]
For fixed $r$ and $m$, the exact error ratio between successive levels is
\begin{equation}
 \frac{\norm{G_{n+1,r}^{(m)}-\Up^{(m)}}_\infty}
 {\norm{G_{n,r}^{(m)}-\Up^{(m)}}_\infty}=4^{-r}
 \label{eq:Gaussian-ratio}
\end{equation}
throughout the stable range.  Thus $r=1,2,3,\ldots$ gives the geometric rates
\[
 4^{-n},\ 16^{-n},\ 64^{-n},\ldots.
\]
\end{corollary}

\subsection{Explicit Gaussian rules}

\begin{example}[One node: ordinary truncation]
For $r=1$, $\pi_1(x)=x$, so $Q_1=\delta_0$ and $h_1=1/9$.  Hence $G_{n,1}=p_n$, and \cref{thm:gaussian-exact} is exactly \cref{thm:ordinary-exact}.
\end{example}

\begin{example}[Two nodes: a variance-matched $16^{-n}$ scheme]
For $r=2$,
\[
 \pi_2(x)=x^2-\frac19,
 \qquad
 Q_2=\frac12\delta_{-1/3}+\frac12\delta_{1/3},
 \qquad
 h_2=\frac{32}{2025}.
\]
Therefore
\begin{equation}
 \boxed{
 G_{n,2}(x)=\frac12p_n\!\left(x-\frac{2^{-n}}3\right)
 +\frac12p_n\!\left(x+\frac{2^{-n}}3\right).}
 \label{eq:Gn2}
\end{equation}
For $n\ge m+5$,
\begin{equation}
 \boxed{
 \norm{G_{n,2}^{(m)}-\Up^{(m)}}_\infty
 =\frac4{6075}\,2^{\binom{m+5}{2}}16^{-n}.}
 \label{eq:Gn2-error}
\end{equation}
In particular,
\begin{equation}
 \norm{G_{n,2}-\Up}_\infty
 =\frac{4096}{6075}\,16^{-n},
 \qquad n\ge5,
 \label{eq:Gn2-uniform}
\end{equation}
with negative equality at $x=-15/16$.
\end{example}

\begin{example}[Three nodes: a $64^{-n}$ scheme]
For $r=3$,
\[
 \pi_3(x)=x^3-\frac{19}{75}x.
\]
The Gaussian nodes and weights are
\begin{equation}
 0,\ \pm\sqrt{\frac{19}{75}},
 \qquad
 \frac{32}{57},\ \frac{25}{114},\ \frac{25}{114},
 \label{eq:G3-nodes-weights}
\end{equation}
respectively, and
\[
 h_3=\frac{19808}{7441875},
 \qquad
 \frac{h_3}{6!}=\frac{1238}{334884375}.
\]
Thus, for $n\ge m+7$,
\begin{equation}
 \boxed{
 \norm{G_{n,3}^{(m)}-\Up^{(m)}}_\infty
 =\frac{1238}{334884375}
 2^{\binom{m+7}{2}}64^{-n}.}
 \label{eq:Gn3-error}
\end{equation}
The extremal point is $-1+2^{-(m+6)}$.
\end{example}

\begin{remark}[Arithmetic form]
The nodes are algebraic roots of rational orthogonal polynomials and the weights lie in the corresponding splitting fields, but the exact norm constants are rational.  This is useful for formalization: one may prove existence and positivity abstractly, while the sharp constant reduces to rational moment and determinant arithmetic.
\end{remark}

\section{Gauss--Lobatto tails: positivity and exact support}
\label{sec:lobatto}

Gaussian nodes lie strictly inside the support, so $G_{n,r}$ has slightly shorter support than $\Up$.  To preserve $[-1,1]$ exactly, force the tail rule to include $-1$ and $1$.  The correct interior nodes are given by orthogonality for the weighted measure $(1-x^2)\dd\mu$.

\subsection{Construction and positivity}

Let $\chi_{r-1}$ be the monic orthogonal polynomial of degree $r-1$ for the measure $\nu$ from \eqref{eq:nu-weighted}.  Its roots
\begin{equation}
 -1<\eta_{r,1}<\cdots<\eta_{r,r-1}<1
 \label{eq:L-interior-nodes}
\end{equation}
are simple.  There is a unique interpolatory rule on the $r+1$ nodes
\begin{equation}
 -1,\ \eta_{r,1},\ldots,\eta_{r,r-1},\ 1
 \label{eq:L-nodes}
\end{equation}
that integrates polynomials of degree at most $r$.  Orthogonality improves its degree of exactness to $2r-1$.  Denote it by
\begin{equation}
 L_r=\lambda_{r,-}\delta_{-1}
 +\sum_{j=1}^{r-1}\lambda_{r,j}\delta_{\eta_{r,j}}
 +\lambda_{r,+}\delta_1.
 \label{eq:Lr}
\end{equation}

\begin{proposition}[Exactness and positive weights]
\label{prop:L-exact-positive}
The rule $L_r$ is exact on polynomials of degree at most $2r-1$, and every weight in \eqref{eq:Lr} is positive.
\end{proposition}

\begin{proof}
Let $p$ have degree at most $2r-1$.  Divide it as
\[
 p(x)=q(x)(x^2-1)\chi_{r-1}(x)+s(x),
 \]
where $\deg q\le r-2$ and $\deg s\le r$.  The first term vanishes at every quadrature node, while
\[
 \int q(x)(x^2-1)\chi_{r-1}(x)\dd\mu(x)
 =-\int q(x)\chi_{r-1}(x)\dd\nu(x)=0.
\]
Thus both the integral and the rule reduce to $s$, proving degree $2r-1$ exactness.

For an interior node $\eta_{r,j}$, the polynomial
\[
 \frac{(1-x^2)\chi_{r-1}(x)^2}{(x-\eta_{r,j})^2}
\]
is nonnegative, has degree $2r-2$, and vanishes at every quadrature node except $\eta_{r,j}$.  Exactness therefore forces the corresponding weight to be positive.  At $1$, use $(1+x)\chi_{r-1}(x)^2$; at $-1$, use $(1-x)\chi_{r-1}(x)^2$.  These have degree $2r-1$, are nonnegative, and isolate the endpoint weights.
\end{proof}

\subsection{The Lobatto Peano kernel}

\begin{proposition}[Positive Lobatto Peano kernel]
\label{prop:L-kernel}
For
\begin{equation}
 \mathcal E_r^L(\phi):=\int\phi\dd L_r-\int\phi\dd\mu,
 \label{eq:L-functional}
\end{equation}
we have $\mathcal E_r^L(\phi)\ge0$ whenever $\phi^{(2r)}\ge0$.  There is a nonnegative kernel $K_r^L$ supported on $[-1,1]$ such that
\begin{equation}
 \mathcal E_r^L(\phi)=\int_{-1}^{1}\phi^{(2r)}(t)K_r^L(t)\dd t,
 \label{eq:L-Peano}
\end{equation}
\begin{equation}
 (K_r^L)^{(2r)}=L_r-\mu,
 \qquad
 \int K_r^L=\frac{\ell_r}{(2r)!}.
 \label{eq:L-kernel-mass}
\end{equation}
\end{proposition}

\begin{proof}
Let $H$ interpolate $\phi$ at the two endpoints and interpolate both $\phi$ and $\phi'$ at the $r-1$ interior nodes.  Then $\deg H\le2r-1$ and
\begin{equation}
 \phi(x)-H(x)
 =\frac{\phi^{(2r)}(\zeta_x)}{(2r)!}
 (x^2-1)\chi_{r-1}(x)^2.
 \label{eq:L-Hermite-remainder}
\end{equation}
The nodal polynomial is nonpositive on $[-1,1]$.  Hence $\phi-H\le0$ for $\phi^{(2r)}\ge0$.  Since $L_r$ is exact on $H$ and uses the same values as $\phi$ at all nodes,
\[
 \mathcal E_r^L(\phi)=\int(H-\phi)\dd\mu\ge0.
\]
For $\phi(x)=x^{2r}$, the difference $x^{2r}-H(x)$ is exactly
\((x^2-1)\chi_{r-1}(x)^2\).  Therefore
\[
 \mathcal E_r^L(x^{2r})
 =\int(1-x^2)\chi_{r-1}(x)^2\dd\mu(x)
 =\ell_r.
\]
The Peano-kernel theorem and the same normalization argument as before give \eqref{eq:L-kernel-mass}.
\end{proof}

Scale $L_r$ and $K_r^L$ by $a=2^{-n}$ as in the Gaussian case, and define
\begin{equation}
 H_{n,r}:=p_n*L_{r,2^{-n}}.
 \label{eq:Hnr}
\end{equation}
We use the letter $H$ to avoid confusing the Lobatto approximant with its kernel.  Because $L_r$ has positive endpoint masses, the support is exactly
\begin{equation}
 \supp(H_{n,r})=[-1,1].
 \label{eq:H-support}
\end{equation}

\begin{theorem}[Sharp support-preserving hierarchy]
\label{thm:L-exact}
Let $r\ge1$, $m\ge0$, and $n\ge m+2r+1$.  Then
\begin{equation}
 \boxed{
 \norm{H_{n,r}^{(m)}-\Up^{(m)}}_\infty
 =\frac{2^{\binom{m+2r+1}{2}}\ell_r}{(2r)!}\,2^{-2rn}.}
 \label{eq:L-exact}
\end{equation}
At $x_{m+2r}=-1+2^{-(m+2r)}$,
\begin{equation}
 \boxed{
 \bigl(H_{n,r}^{(m)}-\Up^{(m)}\bigr)(x_{m+2r})
 =+\frac{2^{\binom{m+2r+1}{2}}\ell_r}{(2r)!}\,2^{-2rn}.}
 \label{eq:L-signed}
\end{equation}
\end{theorem}

\begin{proof}
The scaled Lobatto kernel satisfies
\[
 (K_{r,a}^L)^{(2r)}=L_{r,a}-\mu_a,
 \qquad
 \int K_{r,a}^L=a^{2r}\frac{\ell_r}{(2r)!}.
\]
Therefore
\begin{equation}
 H_{n,r}-\Up=p_n^{(2r)}*K_{r,a}^L.
 \label{eq:L-factorization}
\end{equation}
Differentiate $m$ times and apply the positive-kernel argument and \cref{lem:plateau}.  Equality and sign follow at the center of the saturated plateau.
\end{proof}

\subsection{Explicit support-preserving rules}

\begin{example}[Three-point variance-matched rule]
For $r=2$, $\chi_1(x)=x$.  The nodes are $-1,0,1$ and exactness gives
\begin{equation}
 L_2=\frac1{18}\delta_{-1}+\frac89\delta_0+\frac1{18}\delta_1.
 \label{eq:L2}
\end{equation}
This is the unique symmetric positive law on $\{-1,0,1\}$ with variance $1/9$.  Since
\[
 \ell_2=\mu_2-\mu_4=\frac{56}{675},
 \qquad
 \frac{\ell_2}{4!}=\frac7{2025},
\]
the support-exact finite spline is
\begin{equation}
 \boxed{
 H_{n,2}(x)=\frac1{18}p_n(x+2^{-n})
 +\frac89p_n(x)
 +\frac1{18}p_n(x-2^{-n}).}
 \label{eq:Hn2}
\end{equation}
For $n\ge m+5$,
\begin{equation}
 \boxed{
 \norm{H_{n,2}^{(m)}-\Up^{(m)}}_\infty
 =\frac7{2025}\,2^{\binom{m+5}{2}}16^{-n}.}
 \label{eq:Hn2-error}
\end{equation}
In particular,
\begin{equation}
 \boxed{
 \norm{H_{n,2}-\Up}_\infty
 =\frac{7168}{2025}\,16^{-n},\qquad n\ge5,}
 \label{eq:Hn2-uniform}
\end{equation}
with positive equality at $x=-15/16$.
\end{example}

\begin{example}[Four-point $64^{-n}$ support-exact rule]
For $r=3$,
\[
 \chi_2(x)=x^2-\frac7{75}.
\]
The nodes are
\begin{equation}
 \pm1,\qquad \pm\sqrt{\frac7{75}},
 \label{eq:L3-nodes}
\end{equation}
and their weights are
\begin{equation}
 \frac1{102},\ \frac1{102}
 \quad\text{at }\pm1,
 \qquad
 \frac{25}{51},\ \frac{25}{51}
 \quad\text{at }\pm\sqrt{\frac7{75}}.
 \label{eq:L3-weights}
\end{equation}
Furthermore,
\[
 \ell_3=\frac{78976}{7441875},
 \qquad
 \frac{\ell_3}{6!}=\frac{4936}{334884375}.
\]
Hence, for $n\ge m+7$,
\begin{equation}
 \boxed{
 \norm{H_{n,3}^{(m)}-\Up^{(m)}}_\infty
 =\frac{4936}{334884375}
 2^{\binom{m+7}{2}}64^{-n}.}
 \label{eq:Hn3-error}
\end{equation}
\end{example}

\begin{remark}[Resolution of the variance-matching problem]
The rule \eqref{eq:L2} is positive, supported on the exact normalized tail interval $[-1,1]$, and matches mass, mean, variance, and--by symmetry--the third moment.  After scaling and convolution with $p_n$, it gives the $16^{-n}$ support-exact scheme anticipated but not constructed in the earlier corpus.
\end{remark}

\section{Higher-convex order, sharp weak errors, and scheme comparison}
\label{sec:convex-order}

\subsection{A two-sided quadrature sandwich}

The Gaussian and Lobatto remainder signs imply a clean order relation for observables.

\begin{theorem}[$2r$-convex sandwich]
\label{thm:sandwich}
Let $\phi\in C^{2r}[-1,1]$ with $\phi^{(2r)}\ge0$.  Then
\begin{equation}
 \boxed{
 \int\phi\dd Q_r
 \le \int\phi\dd\mu
 \le \int\phi\dd L_r.}
 \label{eq:sandwich}
\end{equation}
More quantitatively,
\begin{align}
 0\le \int\phi\dd\mu-\int\phi\dd Q_r
 &\le \frac{h_r}{(2r)!}\norm{\phi^{(2r)}}_\infty,
 \label{eq:G-weak-bound}\\
 0\le \int\phi\dd L_r-\int\phi\dd\mu
 &\le \frac{\ell_r}{(2r)!}\norm{\phi^{(2r)}}_\infty.
 \label{eq:L-weak-bound}
\end{align}
Both constants are sharp, with equality for $\phi(x)=x^{2r}$.
\end{theorem}

\begin{proof}
The signs follow from \cref{prop:gauss-kernel,prop:L-kernel}.  Bounding the Peano representations by the kernel masses proves the inequalities.  For $x^{2r}$, the derivative is the constant $(2r)!$ and the errors are $h_r$ and $\ell_r$.
\end{proof}

Let $Z_G\sim Q_r$ and $Z_L\sim L_r$, independent of $S_n$.  Scaling \cref{thm:sandwich} gives, for every $2r$-convex $\phi$,
\begin{equation}
 \boxed{
 \E\phi(S_n+2^{-n}Z_G)
 \le \E\phi(Y)
 \le \E\phi(S_n+2^{-n}Z_L).}
 \label{eq:full-sandwich}
\end{equation}
This is a weak-order statement, not a global pointwise ordering of the three densities.  Pointwise signs depend on the sign of the appropriate derivative of $p_n$; the exact extrema in \cref{thm:gaussian-exact,thm:L-exact} occur where that derivative is saturated and positive.

\subsection{Four representative finite-tail models}

\begin{center}
\begin{tabularx}{0.98\textwidth}{@{}p{0.18\textwidth}p{0.17\textwidth}p{0.18\textwidth}p{0.17\textwidth}X@{}}
\toprule
Scheme & Normalized tail & Matched degrees & Support after convolution & Exact uniform rate and constant\\
\midrule
Ordinary $p_n$ & $\delta_0$ & $0,1$ & radius $1-2^{-n}$ & \(\frac49\,4^{-n}\)\\
Repeated integration $f_n$ & uniform on $[-1,1]$ & $0,1$ & $[-1,1]$ exactly & \(\frac89\,4^{-n}\) (established in corpus)\\
Gaussian $G_{n,2}$ & atoms at $\pm1/3$ & through degree $3$ & strictly inside $[-1,1]$ & \(\frac{4096}{6075}\,16^{-n}\)\\
Lobatto $H_{n,2}$ & weights $1/18,8/9,1/18$ at $-1,0,1$ & through degree $3$ & $[-1,1]$ exactly & \(\frac{7168}{2025}\,16^{-n}\)\\
\bottomrule
\end{tabularx}
\end{center}

Gaussian quadrature has the smaller high-order constant and uses fewer atoms.  Lobatto quadrature pays a constant-factor price for exact endpoint support.  Both improve the geometric ratio from $1/4$ to $1/16$.

\subsection{A design principle}

The constructions expose a general approximation design triangle:
\begin{enumerate}
\item \emph{positivity} is obtained by quadrature weights;
\item \emph{order} is obtained by moment exactness;
\item \emph{support exactness} is obtained by including the endpoints.
\end{enumerate}
For the self-similar Fabius tail, these choices are unusually transparent because the tail at every stage is a scaled copy of one fixed compactly supported measure and because the prefix derivatives have exact Thue--Morse plateaux.

\section{An all-orders expansion for the ordinary prefix}
\label{sec:reciprocal-sinc}

The exact law in \cref{thm:ordinary-exact} comes from a positive physical-space kernel.  The full correction series is more naturally read from the reciprocal of the tail sinc product.

\subsection{Reciprocal-sinc coefficients}

Near the origin, define
\begin{equation}
 \frac1{\Phi(z)}=\sum_{j=0}^{\infty}b_jz^{2j},
 \qquad b_0=1.
 \label{eq:B-reciprocal}
\end{equation}
Euler's product for $\sin z$ gives
\begin{align}
 \log\frac1{\Phi(z)}
 &=-\sum_{m=1}^{\infty}\log\sinc(2^{-m}z)\nonumber\\
 &=\sum_{k=1}^{\infty}\gamma_kz^{2k},
 \label{eq:log-reciprocal-Phi}
\end{align}
where
\begin{equation}
 \boxed{
 \gamma_k=\frac{\zeta(2k)}{k\pi^{2k}(4^k-1)}.}
 \label{eq:gamma-k}
\end{equation}
The ratios $\zeta(2k)/\pi^{2k}$ are rational, so every $\gamma_k$ and every $b_j$ is rational.  Logarithmic differentiation yields the effective recurrence
\begin{equation}
 \boxed{
 b_j=\frac1j\sum_{k=1}^{j}k\gamma_kb_{j-k},
 \qquad j\ge1.}
 \label{eq:b-recurrence}
\end{equation}
The first coefficients are
\begin{align}
 b_0&=1,&
 b_1&=\frac1{18},&
 b_2&=\frac{31}{16200},\label{eq:b-low-1}\\
 b_3&=\frac{2347}{42865200},&
 b_4&=\frac{1904369}{1311675120000}.\label{eq:b-low-2}
\end{align}

\subsection{The expansion theorem}

\begin{theorem}[Reciprocal-sinc expansion]
\label{thm:prefix-expansion}
Fix integers $m\ge0$ and $J\ge1$.  As $n\to\infty$,
\begin{equation}
 \boxed{
 \left\|
 p_n^{(m)}-
 \sum_{j=0}^{J-1}(-1)^jb_j4^{-jn}\Up^{(m+2j)}
 \right\|_\infty
 =\Oh_{m,J}(4^{-Jn}).}
 \label{eq:prefix-expansion}
\end{equation}
The implied constant is independent of $n$.
\end{theorem}

\begin{proof}
Let
\[
 B_{J-1}(z):=\sum_{j=0}^{J-1}b_jz^{2j}.
\]
By the defining Taylor expansion,
\begin{equation}
 1-\Phi(z)B_{J-1}(z)=\Oh(z^{2J})
 \qquad(z\to0).
 \label{eq:local-B-remainder}
\end{equation}
Since $\abs{\Phi(z)}\le1$ on the real axis and $B_{J-1}$ is a fixed polynomial, the local bound extends to the global estimate
\begin{equation}
 \abs{1-\Phi(z)B_{J-1}(z)}\le C_J\abs z^{2J},
 \qquad z\in\R.
 \label{eq:global-B-remainder}
\end{equation}
Indeed, use \eqref{eq:local-B-remainder} for $\abs z\le1$ and the trivial polynomial bound for $\abs z\ge1$.

Write $a=2^{-n}$.  From \eqref{eq:Phi-factor},
\begin{align}
 P_n(\xi)-\Phi(\xi)B_{J-1}(a\xi)
 &=P_n(\xi)\left[1-\Phi(a\xi)B_{J-1}(a\xi)\right].
 \label{eq:Fourier-expansion-remainder}
\end{align}
After multiplying by $\abs\xi^m$, \eqref{eq:global-B-remainder} gives
\begin{equation}
 \abs\xi^m
 \abs{P_n(\xi)-\Phi(\xi)B_{J-1}(a\xi)}
 \le C_Ja^{2J}\abs\xi^{m+2J}\abs{P_n(\xi)}.
 \label{eq:Fourier-majorant}
\end{equation}
Choose a fixed $N>m+2J+1$.  For $n\ge N$,
\[
 \abs{P_n(\xi)}\le\prod_{k=1}^{N}\abs{\sinc(2^{-k}\xi)},
\]
and the right-hand side decays as $\Oh(\abs\xi^{-N})$.  Thus
\(\abs\xi^{m+2J}P_n(\xi)\) has an $L^1$ majorant independent of $n$.  Fourier inversion yields an $\Oh(a^{2J})$ uniform error.  Finally,
\[
 \mathcal F[\Up^{(m+2j)}](\xi)
 =(\ii\xi)^{m+2j}\Phi(\xi)
 =(\ii\xi)^m(-1)^j\xi^{2j}\Phi(\xi),
\]
which gives the signs in \eqref{eq:prefix-expansion}.
\end{proof}

The first four terms are therefore
\begin{equation}
 \boxed{
 \begin{aligned}
 p_n^{(m)}={}&\Up^{(m)}
 -\frac{\Up^{(m+2)}}{18\,4^n}
 +\frac{31\,\Up^{(m+4)}}{16200\,16^n}\\
 &-\frac{2347\,\Up^{(m+6)}}{42865200\,64^n}
 +\Oh_m(256^{-n}).
 \end{aligned}}
 \label{eq:prefix-first-terms}
\end{equation}
The exact law of \cref{thm:ordinary-exact} is compatible with the first term but is stronger: at the saturated dyadic point the physical-space kernel gives a finite-stage equality rather than merely an asymptotic statement.

\subsection{Comparison with the repeated-integration series}

The repeated-integration approximation has a different tail multiplier and begins
\[
 f_n=\Up+\frac1{9}4^{-n}\Up''+\frac{2}{2025}16^{-n}\Up^{(4)}+\cdots.
\]
Thus ordinary truncation and repeated integration approach $\Up$ from opposite second-order directions.  The two-node Gaussian and three-node Lobatto tails cancel that second-order term entirely, leaving fourth-order coefficients
\begin{equation}
 G_{n,2}=\Up-\frac4{6075}16^{-n}\Up^{(4)}+\Oh(64^{-n}),
 \label{eq:G2-leading}
\end{equation}
\begin{equation}
 H_{n,2}=\Up+\frac7{2025}16^{-n}\Up^{(4)}+\Oh(64^{-n}).
 \label{eq:H2-leading}
\end{equation}
The signs agree with the Gaussian/Lobatto Peano kernels and with the exact extrema.

\section{$q=1/4$ Richardson acceleration and Gaussian $q$-binomials}
\label{sec:q-richardson}

The corpus already contains Gaussian $q$-binomial formulas at base $q=1/2$, arising from geometric interpolation in the dyadic spatial scale.  The ordinary-prefix expansion exposes a different geometric variable: its error powers are
\[
 4^{-n},16^{-n},64^{-n},\ldots.
\]
Hence the natural extrapolation base is
\begin{equation}
 q:=\frac14.
 \label{eq:q-quarter}
\end{equation}
This produces a second, conceptually distinct appearance of Gaussian $q$-binomial coefficients.

\subsection{Geometric Lagrange weights}

For $s\ge1$ and $0\le j\le s-1$, let $\lambda_{s,j}$ be the Lagrange cardinal weight for evaluation at $0$ on the nodes $1,q,\ldots,q^{s-1}$:
\begin{equation}
 \lambda_{s,j}:=\prod_{\substack{0\le\ell\le s-1\\\ell\ne j}}
 \frac{-q^\ell}{q^j-q^\ell}.
 \label{eq:lambda-Lagrange}
\end{equation}
A direct factorization gives
\begin{equation}
 \boxed{
 \lambda_{s,j}
 =\frac{(-1)^{s-1-j}q^{\binom{s-j}{2}}}
 {(q;q)_j(q;q)_{s-1-j}}
 =\frac{(-1)^{s-1-j}q^{\binom{s-j}{2}}}{(q;q)_{s-1}}
 \qbinom{s-1}{j}{q}.}
 \label{eq:lambda-qbinomial}
\end{equation}
Here
\((q;q)_k=\prod_{h=1}^{k}(1-q^h)\).

Because these are interpolation weights,
\begin{equation}
 \sum_{j=0}^{s-1}\lambda_{s,j}q^{kj}
 =\begin{cases}
 1,&k=0,\\
 0,&1\le k<s.
 \end{cases}
 \label{eq:lambda-cancellation}
\end{equation}
For the first un-cancelled power,
\begin{equation}
 \sum_{j=0}^{s-1}\lambda_{s,j}q^{sj}
 =(-1)^{s-1}q^{\binom s2}.
 \label{eq:first-uncancelled}
\end{equation}

\begin{proof}[Proof of \eqref{eq:first-uncancelled}]
Interpolate $f(x)=x^s$ at the $s$ nodes $q^j$.  Since $f(0)=0$, the interpolation remainder at $0$ gives
\[
 -\sum_j\lambda_{s,j}q^{sj}
 =\prod_{j=0}^{s-1}(0-q^j)
 =(-1)^sq^{0+1+\cdots+(s-1)}.
\]
Rearranging proves the identity.
\end{proof}

\paragraph{Formal crosswalk.}
Set \(r=s-1\) and \(k=j\).  The compiled module
\path{GeometricQBinomialLagrange.lean} defines a numerator that is zero
outside \(k\le r\) and, inside that range, is exactly
\( (-1)^{r-k}q^{\binom{r-k+1}{2}}
\texttt{gaussianBinomial}\ q\ r\ (r-k)\).
Its \texttt{reversed\_finite\_qBinomial\_theorem} proves the
denominator-free identity with product
\(\prod_{h<r}(z-q^{h+1})\).  If \(k\mapsto q^k\) is injective on
\(\{0,\ldots,r\}\),
\texttt{geometricLagrangeWeight\_eq\_gaussianBinomial\_div} identifies
the normalized numerator with the Lagrange weight, and
\texttt{sum\_geometricLagrangeWeight\_mul\_pow\_eq\_gaussianBinomial}
proves for every \(d>0\)
\[
 \sum_{k=0}^{r}\lambda_{s,k}(q)(q^k)^d
 =(-1)^r q^{r(r+1)/2}
   \texttt{gaussianBinomial}\ q\ (d-1)\ r.
\]
This discharges the finite weight and all positive residual-moment
obligations generically, hence also at \(q=1/4\) whenever the node
injectivity side condition is supplied.  The literal formal coefficient is
reverse-indexed, \texttt{gaussianBinomial} \(q\ r\ (r-k)\); the displayed
\(\qbinom{r}{k}{q}\) in \eqref{eq:lambda-qbinomial} must not be read as an
implicit use of a generic symmetry theorem.  At \(q=1/2\), the existing
\texttt{discreteLimitWeight\_eq\_geometricLagrangeWeight} supplies the
half-base Toeplitz specialization.  None of this proves
\cref{thm:q-richardson}: its analytic expansion, weighted remainder, and
convergence assertions remain frontier.

\subsection{Accelerated combinations}

Define the signed combination
\begin{equation}
 R_{n,s}:=\sum_{j=0}^{s-1}\lambda_{s,j}p_{n+j}.
 \label{eq:Rns}
\end{equation}
Although $R_{n,s}$ is generally not positive, it retains exact mass one because the weights sum to one.

\begin{theorem}[$q$-binomial depth acceleration]
\label{thm:q-richardson}
Fix $m\ge0$ and $s\ge1$.  Then
\begin{equation}
 \boxed{
 R_{n,s}^{(m)}
 =\Up^{(m)}
 -b_sq^{\binom s2}q^{sn}\Up^{(m+2s)}
 +\Oh_{m,s}(q^{(s+1)n}),
 \qquad q=\frac14.}
 \label{eq:q-richardson-asymptotic}
\end{equation}
In particular,
\begin{equation}
 \norm{R_{n,s}^{(m)}-\Up^{(m)}}_\infty=\Oh_{m,s}(4^{-sn}).
 \label{eq:q-richardson-rate}
\end{equation}
\end{theorem}

\begin{proof}
Apply \cref{thm:prefix-expansion} to $p_{n+j}^{(m)}$ through order $s$.  The coefficient of
\(\Up^{(m+2k)}\) is
\[
 (-1)^kb_kq^{kn}\sum_{j=0}^{s-1}\lambda_{s,j}q^{kj}.
\]
The moment identities \eqref{eq:lambda-cancellation} cancel $1\le k<s$.  At $k=s$, use \eqref{eq:first-uncancelled}; the two signs multiply to $-1$.  The uniform remainder remains of order $q^{(s+1)n}$ because $s$ is fixed.
\end{proof}

\subsection{Low-order filters}

For $s=2$,
\begin{equation}
 (\lambda_{2,0},\lambda_{2,1})=\left(-\frac13,\frac43\right),
 \end{equation}
so
\begin{equation}
 \boxed{
 R_{n,2}=\frac{4p_{n+1}-p_n}{3}
 =\Up-\frac{31}{64800}16^{-n}\Up^{(4)}+\Oh(64^{-n}).}
 \label{eq:Rn2}
\end{equation}
For $s=3$,
\begin{equation}
 (\lambda_{3,0},\lambda_{3,1},\lambda_{3,2})
 =\left(\frac1{45},-\frac49,\frac{64}{45}\right),
\end{equation}
so
\begin{equation}
 \boxed{
 R_{n,3}=\frac{p_n-20p_{n+1}+64p_{n+2}}{45}
 =\Up-\frac{2347}{2743372800}64^{-n}\Up^{(6)}
 +\Oh(256^{-n}).}
 \label{eq:Rn3}
\end{equation}
For $s=4$, the weights are
\begin{equation}
 \left(-\frac1{2835},\frac4{135},-\frac{64}{135},\frac{4096}{2835}\right).
 \label{eq:R4-weights}
\end{equation}

\begin{remark}[Two bases, two mechanisms]
The base $1/2$ in the corpus's exact dyadic formulas is the spatial refinement ratio.  The base $1/4$ here is the leading \emph{error} ratio, forced by symmetry and the first nonzero tail moment.  Both yield Gaussian binomial rows, but they act on different variables and solve different problems.
\end{remark}

\section{Exact and high-precision verification}
\label{sec:verification}

The proofs above are analytic.  Exact computation nevertheless provides useful independent checks, especially for signs, powers of two, and index thresholds.

\subsection{Exact rational checks}

At $x=-3/4$, $\Up(x)=F(1/4)=5/72$.  Direct evaluation of the Thue--Morse spline \eqref{eq:TM-spline} gives:

\begin{center}
\begin{tabular}{@{}rrrr@{}}
\toprule
$n$ & $p_n(-3/4)$ & $p_n(-3/4)-\Up(-3/4)$ & predicted value\\
\midrule
3 & $1/16$ & $-1/144$ & $-(4/9)4^{-3}$\\
4 & $13/192$ & $-1/576$ & $-(4/9)4^{-4}$\\
5 & $53/768$ & $-1/2304$ & $-(4/9)4^{-5}$\\
6 & $71/1024$ & $-1/9216$ & $-(4/9)4^{-6}$\\
7 & $853/12288$ & $-1/36864$ & $-(4/9)4^{-7}$\\
8 & $3413/49152$ & $-1/147456$ & $-(4/9)4^{-8}$\\
\bottomrule
\end{tabular}
\end{center}

For the two-node Gaussian rule, equality is attained at $x=-15/16$.  Exact rational arithmetic gives:

\begin{center}
\begin{tabular}{@{}rrr@{}}
\toprule
$n$ & $G_{n,2}(-15/16)-\Up(-15/16)$ & exact prediction\\
\midrule
5 & $-1/1555200$ & $-(4096/6075)16^{-5}$\\
6 & $-1/24883200$ & $-(4096/6075)16^{-6}$\\
7 & $-1/398131200$ & $-(4096/6075)16^{-7}$\\
8 & $-1/6370099200$ & $-(4096/6075)16^{-8}$\\
9 & $-1/101921587200$ & $-(4096/6075)16^{-9}$\\
\bottomrule
\end{tabular}
\end{center}

The support-preserving rule has the opposite sign at the same point:
\begin{equation}
 H_{5,2}(-15/16)-\Up(-15/16)=\frac7{2073600},
 \label{eq:H52-check}
\end{equation}
followed by an exact factor $1/16$ at every subsequent level.

\subsection{Third-order checks}

For $r=3$, the nodes involve square roots.  Eighty-digit evaluation of the algebraic-node mixtures agrees with the exact constants.  For example,
\begin{center}
\begin{tabular}{@{}rcc@{}}
\toprule
$n$ & $\abs{G_{n,3}(-63/64)-\Up(-63/64)}$ & exact value\\
\midrule
7 & $1.7627708126288924\times10^{-12}$ & $619/351151718400000$\\
8 & $2.7543293947326444\times10^{-14}$ & $619/22473709977600000$\\
9 & $4.3036396792697570\times10^{-16}$ & $619/1438317438566400000$\\
\bottomrule
\end{tabular}
\end{center}
The ratio is exactly $1/64$.  The Lobatto errors at the same point are four times larger to first displayed precision and have positive sign, in agreement with \eqref{eq:Hn3-error}.

\subsection{Minimal exact-arithmetic pseudocode}

The following code fragment computes the moment sequence, the Thue--Morse prefix, and the low-order constants using rational arithmetic.

\begin{lstlisting}[style=code,language=Python,caption={Exact arithmetic for moments and dyadic prefix values.}]
from fractions import Fraction
from math import comb, factorial

def tm_sign(j):
    return -1 if j.bit_count() & 1 else 1

def positive_power(x, k):
    return x**k if x > 0 else Fraction(0)

def prefix_density(n, x):
    x = Fraction(x)
    A = Fraction(1) - Fraction(1, 2**n)
    s = Fraction(0)
    for j in range(2**n):
        knot = Fraction(j, 2**(n-1))
        s += tm_sign(j) * positive_power(x + A - knot, n-1)
    return Fraction(2**(n*(n-1)//2), factorial(n-1)) * s

# Given exact cumulants kappa[0..N], moments obey:
def moments_from_cumulants(kappa, N):
    mu = [Fraction(0)] * (N + 1)
    mu[0] = Fraction(1)
    for n in range(1, N + 1):
        mu[n] = sum(Fraction(comb(n-1, k-1))
                    * kappa[k] * mu[n-k]
                    for k in range(1, n + 1))
    return mu

mu2 = Fraction(1, 9)
mu4 = Fraction(19, 675)
mu6 = Fraction(583, 59535)
h2 = mu4 - mu2**2
h3 = mu6 - mu4**2 / mu2
ell2 = mu2 - mu4
assert h2 == Fraction(32, 2025)
assert h3 == Fraction(19808, 7441875)
assert ell2 == Fraction(56, 675)
\end{lstlisting}

\subsection{Verification philosophy}

The numerical checks are not used to infer the theorems.  Their role is narrower:
\begin{enumerate}
\item detect normalization mismatches between $\sin z/z$ and $\sin(\pi z)/(\pi z)$;
\item confirm the signs of the Gaussian and Lobatto kernels;
\item verify the exact dyadic extremal points;
\item catch missing factorials in Peano-kernel masses;
\item test the geometric ratios $4^{-r}$ before formalization.
\end{enumerate}
The actual certificates are the positive-kernel factorizations and the Thue--Morse plateau.

\section{Consequences and additional deductions}
\label{sec:consequences}

\subsection{Exact dyadic values reappear as extremal certificates}

The sharp points are not arbitrary.  The order-$d$ derivative plateau is centered at
\[
 x_d=-1+2^{-d},
\]
which corresponds under $F(x)=\Up(x-1)$ to the dyadic argument $2^{-d}$.  Thus the exact dyadic arithmetic of $F$ provides the values at precisely the points where the finite convolution errors saturate.  For example,
\[
 \Up(-3/4)=F(1/4)=\frac5{72},
 \qquad
 \Up(-15/16)=F(1/16),
 \qquad
 \Up(-63/64)=F(1/64).
\]
The approximation theory and the dyadic-value theory therefore meet at a natural sequence of certificate points.

More generally, the derivative relation for the global Fabius extension implies that higher derivatives of $\Up$ at these points reduce to scaled dyadic values.  The exact-error constants are consequently compatible with the finite derivative tower already developed in the corpus.

\subsection{A rational hierarchy despite algebraic nodes}

The construction separates into three arithmetic layers:
\begin{enumerate}
\item The moments $\mu_{2k}$ are rational because their cumulants contain only Bernoulli numbers and $4^k-1$.
\item The orthogonal polynomials $\pi_r$ and $\chi_{r-1}$ have rational coefficients because their defining Gram systems have rational entries.
\item The nodes and weights are algebraic, but the sharp norm constants $h_r/(2r)!$ and $\ell_r/(2r)!$ are rational determinant ratios.
\end{enumerate}
Thus the approximation itself may use algebraic shifts while its certified worst-case error remains a rational multiple of a power of two.

\begin{corollary}[Rational sharp constants]
For every $r\ge1$ and $m\ge0$, the normalized quantities
\begin{align}
 2^{2rn-\binom{m+2r+1}{2}}
 \norm{G_{n,r}^{(m)}-\Up^{(m)}}_\infty
 &=\frac{h_r}{(2r)!},\label{eq:rational-G-normalized}\\
 2^{2rn-\binom{m+2r+1}{2}}
 \norm{H_{n,r}^{(m)}-\Up^{(m)}}_\infty
 &=\frac{\ell_r}{(2r)!}
 \label{eq:rational-L-normalized}
\end{align}
are positive rational numbers throughout the stable range.
\end{corollary}

\subsection{A compact-support hierarchy of increasingly smooth splines}

For fixed $r$, both $G_{n,r}$ and $H_{n,r}$ are finite mixtures of copies of $p_n$.  Hence they are piecewise polynomials of degree $n-1$ and belong to $C^{n-2}(\R)$.  The number of mixture components is fixed while $n$ grows, so the cost of the tail correction is independent of the spline degree.  Using \eqref{eq:TM-spline}, every component has an explicit finite Thue--Morse representation.  For example,
\begin{align}
 H_{n,2}(x)
 &=\frac{2^{\binom n2}}{(n-1)!}
 \sum_{j=0}^{2^n-1}\TM_j
 \Bigg[
 \frac1{18}\pospart{x+2^{-n}+1-2^{-n}-j2^{1-n}}^{n-1}\nonumber\\
 &\hspace{7em}+\frac89\pospart{x+1-2^{-n}-j2^{1-n}}^{n-1}\nonumber\\
 &\hspace{7em}+\frac1{18}\pospart{x-2^{-n}+1-2^{-n}-j2^{1-n}}^{n-1}
 \Bigg].
 \label{eq:Hn2-TM}
\end{align}
Thus the new positive support-exact approximant is still an explicitly computable Thue--Morse spline.

\subsection{Uniform approximation of integrals and observables}

Let $X_n^G=S_n+2^{-n}Z_G$ and $X_n^L=S_n+2^{-n}Z_L$.  For an arbitrary $\phi\in C^{2r}(\R)$, the Peano kernels yield
\begin{align}
 \abs{\E\phi(Y)-\E\phi(X_n^G)}
 &\le 2^{-2rn}\frac{h_r}{(2r)!}\norm{\phi^{(2r)}}_\infty,
 \label{eq:weak-G-full}\\
 \abs{\E\phi(X_n^L)-\E\phi(Y)}
 &\le 2^{-2rn}\frac{\ell_r}{(2r)!}\norm{\phi^{(2r)}}_\infty.
 \label{eq:weak-L-full}
\end{align}
These bounds require no derivative plateau and therefore hold at every $n$.  The density norm theorems are stronger and more specialized: they exploit the exact shape of $p_n^{(m+2r)}$ to prove that the generic Peano constant is attained.

\subsection{Asymptotic bracketing of smooth observables}

For a $2r$-convex observable, \eqref{eq:full-sandwich} gives a certified bracket at every level.  In the first nontrivial case $r=2$,
\begin{align}
 \frac12\E\left[
 \phi\!\left(S_n-\frac{2^{-n}}3\right)
 +\phi\!\left(S_n+\frac{2^{-n}}3\right)
 \right]
 &\le \E\phi(Y)\nonumber\\
 &\le
 \frac1{18}\E\phi(S_n-2^{-n})
 +\frac89\E\phi(S_n)
 +\frac1{18}\E\phi(S_n+2^{-n})
 \label{eq:r2-observable-bracket}
\end{align}
whenever $\phi^{(4)}\ge0$.  The bracket width is at most
\begin{equation}
 \left(\frac4{6075}+\frac7{2025}\right)16^{-n}
 \norm{\phi^{(4)}}_\infty
 =\frac5{1215}16^{-n}\norm{\phi^{(4)}}_\infty.
 \label{eq:r2-bracket-width}
\end{equation}
This gives a practical positive-probability certificate without evaluating the full infinite convolution.

\subsection{Why $q=1/4$ is unavoidable}

The law $\mu$ is centered, so every symmetric tail surrogate automatically matches degree one.  The first possible defect is the variance and scales as $a^2=4^{-n}$.  Once the variance is matched, the next possible defect is the fourth moment and scales as $a^4=16^{-n}$.  Thus $q=1/4$ is not merely a convenient extrapolation parameter: it is the square of the dyadic spatial scale and the fundamental ratio of the even-moment hierarchy.  The Richardson weights in \eqref{eq:lambda-qbinomial} are therefore the natural Gaussian-binomial coefficients for symmetric tail errors.

\section{Lean formalization roadmap}
\label{sec:formalization}

The proofs were chosen to decompose into reusable lemmas that fit the existing formal development.  A plausible formalization order is as follows.

\subsection{Phase I: probability measures and scaling}

\begin{enumerate}
\item Define the up-law $\mu$ as the distribution of the established random series and define its scaled pushforward $\mu_a$.
\item Prove the exact tail split
\[
 \mu=\mathcal L(S_n)*\mu_{2^{-n}}.
\]
\item Connect the density of $S_n$ to the existing finite box-spline object $p_n$.
\end{enumerate}

Most of this layer should be available from the current infinite-convolution and probability modules.

\subsection{Phase II: the ordinary-prefix potential}

\begin{enumerate}
\item Define
\[
 J(t)=\frac12\int(\abs{t-y}-\abs t)\dd\mu(y).
\]
\item Prove $J\ge0$ by Jensen, compact support by the support of $\mu$, and evenness by symmetry.
\item Formalize $J''=\mu-\delta_0$ either in distributions or through a twice-integrated identity against compactly supported test functions.
\item Prove the mass identity $\int J=\mu_2/2$.
\item Combine the existing derivative plateau with a nonnegative-kernel convolution lemma to obtain \cref{thm:ordinary-exact}.
\end{enumerate}

The only delicate issue is the a.e. nature of the finite spline derivative at knots.  One clean route is to state the plateau for the $L^\infty$ representative and then use continuity of the final convolution.

\subsection{Phase III: abstract positive quadrature}

A reusable finite-dimensional module can avoid committing immediately to explicit roots.

\begin{enumerate}
\item Develop monic orthogonal polynomials for a compactly supported probability measure with positive moments.
\item Prove the Gram determinant formula $h_r=\Delta_{r+1}/\Delta_r$.
\item Import or prove Gaussian quadrature existence, positivity, and exactness.
\item Formalize Hermite interpolation and the sign of the Gaussian remainder.
\item Define the Peano kernel by the truncated-power formula and prove nonnegativity and mass.
\item Repeat with the weighted measure $(1-x^2)\mu$ for Lobatto quadrature, including the positive-weight proof from isolating nonnegative polynomials.
\end{enumerate}

The density error theorem then becomes a generic lemma:

\frontierbox{If a signed compactly supported measure $\sigma$ has a nonnegative $2r$th-order potential $K$ of mass $M$, and if a density $p$ has a derivative plateau of height $C$ covering $\supp K$, then
\[
 \norm{p*\sigma}_\infty=CM.
\]
}

This abstraction would also subsume the existing repeated-integration theorem and the ordinary-prefix theorem.

\subsection{Phase IV: exact low-order certificates}

The cases $r=2,3$ require only rational moment arithmetic plus elementary algebra:
\begin{enumerate}
\item verify the moments in \eqref{eq:moments-low};
\item prove the displayed orthogonal polynomials;
\item verify nodes and weights by moment exactness;
\item normalize the rational constants;
\item instantiate the generic plateau theorem.
\end{enumerate}
These cases can be formalized before a fully general quadrature library is complete.

\subsection{Phase V: reciprocal sinc and $q$-binomial acceleration}

\begin{enumerate}
\item Define the formal reciprocal coefficients $b_j$ from the power series of $1/\Phi$.
\item Prove the logarithmic coefficient formula \eqref{eq:gamma-k} from Euler's sine product.
\item Establish the fixed-order Fourier majorant used in \cref{thm:prefix-expansion}.
\item Use the compiled \path{GeometricQBinomialLagrange.lean} bridge for the geometric Lagrange weights and all positive residual moments.
\item Combine the two finite expansions to prove \cref{thm:q-richardson}.
\end{enumerate}
The finite algebraic step is now available at arbitrary ratio under node
injectivity, including \(q=1/4\); the remaining step is the analytic
combination with controlled remainders.  Its \(q=1/2\) specialization is the
existing half-base Toeplitz bridge.

\section{Open problems and next frontier}
\label{sec:open}

The present results close two gaps but open a larger approximation theory for self-similar smooth laws.

\subsection{Arithmetic of the Hankel constants}

The sequences
\begin{equation}
 h_r=\frac{\Delta_{r+1}}{\Delta_r},
 \qquad
 \ell_r=\frac{\Lambda_r}{\Lambda_{r-1}}
 \label{eq:two-constant-sequences}
\end{equation}
are positive rationals built from Bernoulli numbers and Mersenne factors.  Their reduced denominators, odd-prime valuations, and possible normalized integrality properties appear unexplored.  Natural questions include:
\begin{enumerate}
\item Is there a useful factorization of $\Delta_r$ or $\Lambda_r$ into predictable Mersenne components?
\item Do suitable powers of two and products $\prod_{j\le r}(4^j-1)$ clear all denominators sharply?
\item Are the numerators related to known moment or continued-fraction sequences?
\end{enumerate}
A Jacobi continued fraction for $M(t)$ or for the Cauchy transform of $\mu$ could encode these determinants more efficiently.

\subsection{Asymptotics as the quadrature order grows}

The fixed-$r$, large-$n$ theory is exact.  A different regime lets $r\to\infty$.  The up density is positive in the interior but flat at the endpoints, so the zero distribution of $\pi_r$ should be governed at first order by logarithmic potential theory on $[-1,1]$, while endpoint spacing and Christoffel functions may feel the extreme flatness of the weight.  Quantitative questions include:
\begin{enumerate}
\item asymptotics of $h_r$ and $\ell_r$;
\item the smallest distance of Gaussian nodes from $\pm1$;
\item comparison of Gaussian and Lobatto constants;
\item uniform algorithms when $r=r(n)$ grows with the convolution depth.
\end{enumerate}

\subsection{Nested positive tail rules}

Gaussian rules are not generally nested.  Lobatto rules share endpoints but not necessarily interior nodes.  A Fabius-specific Gauss--Kronrod or Patterson-type construction could provide nested positive tail laws, allowing one to reuse shifted prefix evaluations.  The relevant optimization problem is:

\frontierbox{Construct nested symmetric positive measures $Q^{(1)}\subset Q^{(2)}\subset\cdots$ on $[-1,1]$ that maximize moment exactness for the up-law, while keeping exact rational error certificates whenever possible.}

The strong dyadic arithmetic of the moments may make this more tractable than for a generic weight.

\subsection{Positive acceleration versus signed Richardson filters}

The $q$-binomial Richardson rule reaches order $4^{-sn}$ with a signed linear combination of different depths.  Gaussian/Lobatto rules reach the same order using positive mixtures at one depth.  It would be valuable to compare:
\begin{enumerate}
\item total variation of the signed Richardson coefficients;
\item number and location of positive quadrature atoms;
\item exact support;
\item numerical stability near flat endpoints;
\item compatibility with exact rational evaluation.
\end{enumerate}
A hybrid method could use a low-order positive tail correction followed by a short $q=1/4$ extrapolation, reducing both the atom count and the signed coefficient growth.

\subsection{Fourier-decay interaction}

The recent Fourier-decay reports analyze
\[
 \Phi(\xi)=\prod_{k\ge1}\sinc(2^{-k}\xi)
\]
along typical and exceptional rays.  The corrected approximants have transforms
\begin{align}
 \widehat G_{n,r}(\xi)&=P_n(\xi)\widehat Q_r(2^{-n}\xi),\label{eq:G-Fourier}\\
 \widehat H_{n,r}(\xi)&=P_n(\xi)\widehat L_r(2^{-n}\xi).
 \label{eq:H-Fourier}
\end{align}
Their low-frequency errors are controlled by moment matching, while their high-frequency behavior is inherited from the finite product $P_n$.  A two-parameter analysis in $n$ and $\xi$ could determine the largest frequency window on which the relative error remains small.  This is especially relevant because $\Phi$ has many exact zeros, so global relative error requires zero-aware normalization.

\subsection{Optimal support-exact positive rules}

The Lobatto rule is canonical among interpolatory rules with both endpoints fixed and degree $2r-1$.  Other positive support-exact tail measures with more atoms may reduce the sharp $C^m$ constant while keeping the same order.  One can pose a finite moment problem:
\begin{equation}
 \min_Q \abs{\mu_{2r}-Q[x^{2r}]}
 \label{eq:moment-optimization}
\end{equation}
subject to positivity, support in $[-1,1]$, inclusion of the endpoints, and matching moments through degree $2r-1$.  Understanding the extreme points of this feasible set may yield better support-exact constants than Lobatto at the cost of additional atoms.

\subsection{Beyond the dyadic scale}

The method is not intrinsically restricted to base two.  For a self-similar random series
\[
 Y_b=\sum_{k\ge1}b^{-k}U_k,
 \qquad b>1,
\]
when the support geometry and finite-prefix derivative plateaux are controlled, moment-matched quadrature predicts rates $b^{-2rn}$.  What is special about $b=2$ is the exact unique-subset-sum structure, which turns the finite spline signs into Thue--Morse signs and gives the saturated derivative constants in closed form.  A classification of bases or weight sequences admitting an equally sharp plateau theorem would generalize the entire framework.

\section{Conclusion}
\label{sec:conclusion}

The finite dyadic prefix, the infinite sinc product, and the Thue--Morse sequence are three descriptions of one object.  The new approximation theory developed here makes that unity explicit.

The omitted self-similar tail has a nonnegative second-order potential of mass $1/18$.  Convolving this potential against a saturated Thue--Morse derivative plateau gives the exact ordinary-prefix law
\[
 \norm{p_n^{(m)}-\Up^{(m)}}_\infty
 =\frac{2^{\binom{m+3}{2}}}{18}4^{-n}.
\]
Replacing the omitted tail by Gaussian quadrature raises the order from $4^{-n}$ to $4^{-rn}$ while preserving positivity.  Replacing it by Gauss--Lobatto quadrature preserves both positivity and the exact support.  The sharp constants are rational orthogonal-polynomial norms and are attained at explicit dyadic points because the finite prefix has exact derivative plateaux.

On the Fourier side, the reciprocal sinc product generates the complete correction series.  Its powers are geometric in $q=1/4$, and geometric Lagrange interpolation produces Gaussian $q$-binomial Richardson weights.  Hence the four central themes of the corpus meet in one coherent mechanism:
\begin{equation}
 \boxed{
 \begin{gathered}
 \text{dyadic random tail}
 \xrightarrow{\text{moment quadrature}}
 \text{positive high-order spline},\\
 \text{infinite sinc product}
 \xrightarrow{\text{reciprocal series}}
 \text{even-power error expansion},\\
 \text{Thue--Morse prefix sums}
 \xrightarrow{\text{plateau saturation}}
 \text{exact }C^m\text{ constants},\\
 q=\tfrac14\text{ geometric interpolation}
 \xrightarrow{\text{Gaussian binomials}}
 \text{depth acceleration}.
 \end{gathered}}
 \label{eq:four-way-summary}
\end{equation}

The most immediate next step is formalization of the ordinary-prefix potential and the explicit $r=2$ Gaussian and Lobatto schemes.  These cases require little new infrastructure, resolve a documented frontier gap, and provide compact theorem statements with exact rational certificates.

\appendix

\section{Coefficient and constant tables}
\label{app:tables}

\subsection{Moments and cumulants}

\begin{longtable}{@{}rlll@{}}
\toprule
$m$ & $\kappa_{2m}$ & $\mu_{2m}$ & comment\\
\midrule
\endfirsthead
\toprule
$m$ & $\kappa_{2m}$ & $\mu_{2m}$ & comment\\
\midrule
\endhead
1 & $\frac19$ & $\frac19$ & variance\\
2 & $-\frac2{225}$ & $\frac{19}{675}$ & first non-Gaussian correction\\
3 & $\frac{16}{3969}$ & $\frac{583}{59535}$ & used by the three-node Gaussian rule\\
4 & $-\frac{16}{3825}$ & $\frac{132809}{32531625}$ & determines degree-eight data\\
5 & $\frac{256}{33759}$ & $\frac{46840699}{24405225075}$ & rational Mersenne structure\\
6 & $-\frac{353792}{16769025}$ & $\frac{4068990560161}{4133856862760625}$ & contains the Bernoulli numerator $691$\\
\bottomrule
\end{longtable}

The cumulant entries follow directly from \eqref{eq:cumulants}; the displayed higher fractions are included mainly as exact-arithmetic checkpoints.

\subsection{Reciprocal-sinc coefficients}

\begin{center}
\begin{tabular}{@{}rll@{}}
\toprule
$j$ & $\gamma_j$ in $\log(1/\Phi)$ & $b_j$ in $1/\Phi$\\
\midrule
0 & --- & $1$\\
1 & $1/18$ & $1/18$\\
2 & $1/2700$ & $31/16200$\\
3 & $1/178605$ & $2347/42865200$\\
4 & $1/9639000$ & $1904369/1311675120000$\\
5 & $1/478533825$ & $3309760579/88561680752160000$\\
\bottomrule
\end{tabular}
\end{center}

\subsection{Quadrature constants}

\begin{center}
\begin{tabularx}{0.98\textwidth}{@{}rXXll@{}}
\toprule
$r$ & Gaussian polynomial $\pi_r$ & Lobatto interior polynomial $\chi_{r-1}$ & $h_r/(2r)!$ & $\ell_r/(2r)!$\\
\midrule
1 & $x$ & $1$ & $1/18$ & $4/9$\\
2 & $x^2-1/9$ & $x$ & $4/6075$ & $7/2025$\\
3 & $x^3-(19/75)x$ & $x^2-7/75$ & $1238/334884375$ & $4936/334884375$\\
\bottomrule
\end{tabularx}
\end{center}

For $r=1$, the Lobatto convention has the two endpoint rule, exact through degree one.  It is not the ordinary prefix; it corresponds to a first-order support-exact correction.  The higher-order statements in the body begin with the practically important $r=2$ case.

\section{Derivation of the low-order quadrature rules}
\label{app:low-rules}

\subsection{Gaussian order two}

Symmetry gives $\pi_2=x^2-c$.  Orthogonality to $1$ gives $c=\mu_2=1/9$.  The roots are $\pm1/3$, and symmetry plus total mass forces weights $1/2$.  The norm is
\[
 h_2=\mu_4-\mu_2^2
 =\frac{19}{675}-\frac1{81}
 =\frac{32}{2025}.
\]

\subsection{Gaussian order three}

Write $\pi_3=x^3-cx$.  Orthogonality to $x$ gives
\[
 c=\frac{\mu_4}{\mu_2}=\frac{19}{75}.
\]
If the two outer weights are $w$ and the central weight is $1-2w$, moment matching gives
\[
 2w\frac{19}{75}=\frac19,
 \qquad
 w=\frac{25}{114},
 \qquad
 1-2w=\frac{32}{57}.
\]
Finally,
\[
 h_3=\mu_6-\frac{\mu_4^2}{\mu_2}
 =\frac{19808}{7441875}.
\]

\subsection{Lobatto order two}

The interior polynomial is $\chi_1=x$.  Symmetric weights $(a,b,a)$ at $(-1,0,1)$ obey
\[
 2a+b=1,
 \qquad
 2a=\mu_2=\frac19.
\]
Thus $a=1/18$, $b=8/9$, and
\[
 \ell_2=\int(1-x^2)x^2\dd\mu
 =\mu_2-\mu_4
 =\frac{56}{675}.
\]

\subsection{Lobatto order three}

Write $\chi_2=x^2-c$.  Weighted orthogonality to $1$ gives
\[
 0=\int(1-x^2)(x^2-c)\dd\mu
 =\mu_2-\mu_4-c(1-\mu_2),
\]
so $c=7/75$.  Let $a$ be each endpoint weight and $b$ each interior weight.  Then
\[
 a+b=\frac12,
 \qquad
 2a+2b\frac7{75}=\frac19,
\]
which gives $a=1/102$ and $b=25/51$.  Direct substitution yields
\[
 \ell_3=\int(1-x^2)\left(x^2-\frac7{75}\right)^2\dd\mu
 =\frac{78976}{7441875}.
\]

\section{A reusable abstract sharpness lemma}
\label{app:abstract-lemma}

The ordinary, Gaussian, Lobatto, and repeated-integration laws all fit one template.

\begin{theorem}[Positive-kernel plateau principle]
\label{thm:abstract-plateau}
Let $K\in L^1(\R)$ satisfy $K\ge0$, $\supp K\subseteq[-a,a]$, and $\int K=M$.  Let $g$ be essentially bounded and suppose
\[
 \norm g_\infty=C,
 \qquad
 g=C\quad\text{a.e. on }[x_0-a,x_0+a].
\]
Then
\begin{equation}
 \norm{g*K}_\infty=CM,
 \qquad
 (g*K)(x_0)=CM.
 \label{eq:abstract-plateau}
\end{equation}
If an approximation error can be written as $\pm g*K$, its uniform norm and a signed extremum are therefore exact.
\end{theorem}

\begin{proof}
Young's inequality gives $\norm{g*K}_\infty\le CM$.  At $x_0$, the function $g(x_0-t)$ equals $C$ for almost every $t$ in the support of $K$, so the convolution equals $C\int K=CM$.
\end{proof}

In the applications,
\[
 g=p_n^{(m+2r)},
 \qquad
 C=2^{\binom{m+2r+1}{2}},
 \qquad
 a=2^{-n},
\]
and $K$ is a scaled Peano kernel.

\section{Corpus coverage map}
\label{app:coverage}

The report's novelty comparison used the following thematic map.  It is included to make explicit what was treated as established and therefore not redeveloped at length.

\begin{longtable}{@{}p{0.27\textwidth}p{0.66\textwidth}@{}}
\toprule
Corpus component & Material treated as established baseline\\
\midrule
\endfirsthead
\toprule
Corpus component & Material treated as established baseline\\
\midrule
\endhead
Core Fabius/up exposition & Existence, uniqueness, smoothness, support, symmetry, probability representation, moments, Fourier product, Poisson summation, exact dyadic values, $q$-binomial formulas, and corrected endpoint asymptotics.\\
Lean walkthrough & Trusted theorem boundary, module dependencies, executable exact evaluator, and proof architecture.\\
Glossary & Terminology normalization across the entire project.\\
Dyadic formula and representation papers & Highest-bit recursion, global signed extension, coefficient formulas, binary telescopes, and alternative series.\\
Dyadic asymptotic bridge & Exact coefficient extraction, fixed-numerator asymptotics, Lambert coordinate, periodic fluctuations, and Bromwich collapse.\\
Repeated integration & Exact finite box spline with duplicated smallest width, positive second-order Peano kernel, sharp $8/(9\,4^n)$ law, derivative analogues, and its own all-orders expansion.\\
Thue--Morse convergence and formula atlases & Prefix summation, Prouhet cancellation, product/generating formulas, finite differences, parity identities, and q-calculus.\\
$q$-special-functions report & q-Pochhammer symbols, Gaussian binomials, basic hypergeometric structure, and their role in exact Fabius arithmetic.\\
Fourier-decay reports & Exact dyadic factorization, typical-ray exponent, fluctuations, exceptional orbits, covariance, and law-of-the-iterated-logarithm behavior of the sinc product.\\
Non-formalized frontier volume & Consolidated speculative results and explicit gap register; used to identify the ordinary-prefix and moment-matched-tail questions.\\
Archive & Historical variants, superseded drafts, and provenance; duplicate statements were not counted twice.\\
\bottomrule
\end{longtable}

\section{Notation index}
\label{app:notation}

\begin{center}
\begin{tabularx}{0.98\textwidth}{@{}p{0.17\textwidth}X@{}}
\toprule
Symbol & Meaning\\
\midrule
$\Up$ & Rvachev's even compactly supported up-function, normalized as a probability density.\\
$F$ & Bounded Fabius function on $[0,1]$, with $F(x)=\Up(x-1)$.\\
$\mu$ & Probability law with density $\Up$.\\
$Y$ & Random series $\sum_{k\ge1}2^{-k}U_k$ with law $\mu$.\\
$S_n$ & Finite prefix $\sum_{k=1}^n2^{-k}U_k$.\\
$p_n$ & Density of $S_n$, a finite dyadic box spline.\\
$\Phi$ & Infinite sinc product $\widehat\Up$.\\
$P_n$ & Finite sinc product $\widehat p_n$.\\
$\TM_j$ & Signed Thue--Morse value $(-1)^{s_2(j)}$.\\
$C_r$ & Exact derivative norm $2^{\binom{r+1}{2}}$.\\
$x_r$ & Plateau center $-1+2^{-r}$.\\
$J$ & Positive second-order potential for $\mu-\delta_0$.\\
$\pi_r,h_r$ & Monic orthogonal polynomial for $\mu$ and its squared norm.\\
$Q_r,G_{n,r}$ & Gaussian tail rule and corrected finite spline.\\
$\chi_{r-1},\ell_r$ & Monic orthogonal polynomial for $(1-x^2)\mu$ and its squared weighted norm.\\
$L_r,H_{n,r}$ & Gauss--Lobatto tail rule and support-exact corrected spline.\\
$b_j$ & Coefficients of the reciprocal sinc product $1/\Phi$.\\
$R_{n,s}$ & $q=1/4$ Richardson combination of consecutive prefix depths.\\
\bottomrule
\end{tabularx}
\end{center}

\begin{thebibliography}{99}

\bibitem{Fabius1966}
J.~Fabius,
\newblock A probabilistic example of a nowhere analytic $C^\infty$-function,
\newblock \emph{Zeitschrift f\"ur Wahrscheinlichkeitstheorie und Verwandte Gebiete}
\textbf{5} (1966), 173--174.
\newblock DOI: \nolinkurl{10.1007/BF00536652}.

\bibitem{Rvachev1990}
V.~A.~Rvachev,
\newblock Compactly supported solutions of functional-differential equations and their applications,
\newblock \emph{Russian Mathematical Surveys} \textbf{45} (1990), no.~1, 87--120.

\bibitem{Arias1982}
J.~Arias de Reyna,
\newblock Definici\'on y estudio de una funci\'on indefinidamente diferenciable de soporte compacto,
\newblock \emph{Revista de la Real Academia de Ciencias Exactas, F\'isicas y Naturales de Madrid}
\textbf{76} (1982), no.~1, 21--38.
\newblock English translation: \emph{An infinitely differentiable function with compact support: Definition and properties}, arXiv:\nolinkurl{1702.05442}.

\bibitem{Arias2018}
J.~Arias de Reyna,
\newblock Arithmetic of the Fabius function,
\newblock \emph{INTEGERS} \textbf{18} (2018), Article A51;
 arXiv:\nolinkurl{1702.06487}.

\bibitem{AlloucheShallit}
J.-P.~Allouche and J.~Shallit,
\newblock \emph{Automatic Sequences: Theory, Applications, Generalizations},
\newblock Cambridge University Press, 2003.

\bibitem{Gautschi}
W.~Gautschi,
\newblock \emph{Orthogonal Polynomials: Computation and Approximation},
\newblock Numerical Mathematics and Scientific Computation, Oxford University Press, 2004.

\bibitem{DavisRabinowitz}
P.~J.~Davis and P.~Rabinowitz,
\newblock \emph{Methods of Numerical Integration}, second edition,
\newblock Academic Press, 1984.

\bibitem{BrassPetras}
H.~Brass and K.~Petras,
\newblock \emph{Quadrature Theory: The Theory of Numerical Integration on a Compact Interval},
\newblock Mathematical Surveys and Monographs 178, American Mathematical Society, 2011.

\bibitem{Sard}
A.~Sard,
\newblock \emph{Linear Approximation},
\newblock Mathematical Surveys 9, American Mathematical Society, 1963.

\bibitem{Petras2003}
K.~Petras,
\newblock Error bounds for Gaussian and related quadrature and applications to $R$-convex functions,
\newblock \emph{SIAM Journal on Numerical Analysis} \textbf{41} (2003), 2095--2110.

\bibitem{ProveIt}
V.~Reshetnikov,
\newblock \emph{ProveIt: Analysis/FabiusFunction},
\newblock Lean~4 formal development and research corpus, repository snapshot accessed 27 August 2026.
\newblock \url{https://github.com/VladimirReshetnikov/ProveIt/tree/main/Analysis/FabiusFunction}.

\bibitem{MasterDoc}
V.~Reshetnikov,
\newblock \emph{The Fabius Function and Rvachev's Up Function: Self-similarity, probability, exact dyadic arithmetic, and the corrected small-argument asymptotic},
\newblock manuscript in the \texttt{ProveIt} repository, 2026.

\bibitem{FrontierVolume}
V.~Reshetnikov,
\newblock \emph{Non-formalized research frontiers for the Fabius--Rvachev development},
\newblock consolidated manuscript in the \texttt{ProveIt} repository, 2026.

\bibitem{RepeatedIntegration}
V.~Reshetnikov,
\newblock \emph{Repeated Integration and Rvachev's Up-Function: Exact finite splines, positive Peano kernels, and sharp convergence},
\newblock manuscript in the \texttt{ProveIt} repository, 2026.

\bibitem{Haugland}
J.~K.~Haugland,
\newblock Evaluating the Fabius function,
\newblock arXiv:\nolinkurl{1609.07999}, 2016.

\bibitem{Andrews}
G.~E.~Andrews,
\newblock \emph{The Theory of Partitions},
\newblock Encyclopedia of Mathematics and its Applications 2, Addison--Wesley, 1976.

\bibitem{GasperRahman}
G.~Gasper and M.~Rahman,
\newblock \emph{Basic Hypergeometric Series}, second edition,
\newblock Encyclopedia of Mathematics and its Applications 96, Cambridge University Press, 2004.

\end{thebibliography}

\end{document}
